%% Generated by Sphinx.
\def\sphinxdocclass{report}
\documentclass[letterpaper,10pt,english]{sphinxmanual}
\ifdefined\pdfpxdimen
   \let\sphinxpxdimen\pdfpxdimen\else\newdimen\sphinxpxdimen
\fi \sphinxpxdimen=.75bp\relax
\ifdefined\pdfimageresolution
    \pdfimageresolution= \numexpr \dimexpr1in\relax/\sphinxpxdimen\relax
\fi
%% let collapsible pdf bookmarks panel have high depth per default
\PassOptionsToPackage{bookmarksdepth=5}{hyperref}

\PassOptionsToPackage{warn}{textcomp}
\usepackage[utf8]{inputenc}
\ifdefined\DeclareUnicodeCharacter
% support both utf8 and utf8x syntaxes
  \ifdefined\DeclareUnicodeCharacterAsOptional
    \def\sphinxDUC#1{\DeclareUnicodeCharacter{"#1}}
  \else
    \let\sphinxDUC\DeclareUnicodeCharacter
  \fi
  \sphinxDUC{00A0}{\nobreakspace}
  \sphinxDUC{2500}{\sphinxunichar{2500}}
  \sphinxDUC{2502}{\sphinxunichar{2502}}
  \sphinxDUC{2514}{\sphinxunichar{2514}}
  \sphinxDUC{251C}{\sphinxunichar{251C}}
  \sphinxDUC{2572}{\textbackslash}
\fi
\usepackage{cmap}
\usepackage[T1]{fontenc}
\usepackage{amsmath,amssymb,amstext}
\usepackage{babel}



\usepackage{tgtermes}
\usepackage{tgheros}
\renewcommand{\ttdefault}{txtt}



\usepackage[Bjarne]{fncychap}
\usepackage{sphinx}

\fvset{fontsize=auto}
\usepackage{geometry}


% Include hyperref last.
\usepackage{hyperref}
% Fix anchor placement for figures with captions.
\usepackage{hypcap}% it must be loaded after hyperref.
% Set up styles of URL: it should be placed after hyperref.
\urlstyle{same}

\addto\captionsenglish{\renewcommand{\contentsname}{Contents:}}

\usepackage{sphinxmessages}
\setcounter{tocdepth}{1}



\title{target\_setter}
\date{Feb 20, 2026}
\release{1.0}
\author{Chheng Lydiya}
\newcommand{\sphinxlogo}{\vbox{}}
\renewcommand{\releasename}{Release}
\makeindex
\begin{document}

\pagestyle{empty}
\sphinxmaketitle
\pagestyle{plain}
\sphinxtableofcontents
\pagestyle{normal}
\phantomsection\label{\detokenize{index::doc}}



\chapter{Introduction}
\label{\detokenize{introduction:introduction}}\label{\detokenize{introduction::doc}}

\section{Outline}
\label{\detokenize{introduction:outline}}\begin{itemize}
\item {} 
\sphinxAtStartPar
Project background and motivation

\item {} 
\sphinxAtStartPar
System overview and main features

\item {} 
\sphinxAtStartPar
Intended users and operating context

\item {} 
\sphinxAtStartPar
Document organization

\end{itemize}


\chapter{Hardware Specification}
\label{\detokenize{hardware_specification:hardware-specification}}\label{\detokenize{hardware_specification::doc}}
\sphinxAtStartPar
Describe the robot platform hardware, sensors, compute, and network setup.
Populate this section with the exact components used in the system.


\section{Topics to cover}
\label{\detokenize{hardware_specification:topics-to-cover}}\begin{itemize}
\item {} 
\sphinxAtStartPar
Robot base and drive type

\item {} 
\sphinxAtStartPar
Sensors (IMU, encoders, lidar, cameras)

\item {} 
\sphinxAtStartPar
Compute hardware and OS

\item {} 
\sphinxAtStartPar
Networking (Wi\sphinxhyphen{}Fi, Ethernet, IP configuration)

\item {} 
\sphinxAtStartPar
Power system and safety constraints

\end{itemize}


\chapter{Installation and Configuration}
\label{\detokenize{installation_and_configuration:installation-and-configuration}}\label{\detokenize{installation_and_configuration::doc}}
\sphinxAtStartPar
Add step\sphinxhyphen{}by\sphinxhyphen{}step setup instructions for building and running the project.
Include dependencies, workspace setup, and configuration parameters.


\section{Topics to cover}
\label{\detokenize{installation_and_configuration:topics-to-cover}}\begin{itemize}
\item {} 
\sphinxAtStartPar
Prerequisites and dependencies

\item {} 
\sphinxAtStartPar
Workspace setup and build steps

\item {} 
\sphinxAtStartPar
Configuration files and parameters

\item {} 
\sphinxAtStartPar
Running the app and robot nodes

\item {} 
\sphinxAtStartPar
Common setup pitfalls

\end{itemize}


\chapter{Technical Reference}
\label{\detokenize{technical_reference:technical-reference}}\label{\detokenize{technical_reference::doc}}

\section{Architecture}
\label{\detokenize{technical_reference:architecture}}
\noindent\sphinxincludegraphics{{system_overview}.png}


\subsection{Target\_Setter app}
\label{\detokenize{technical_reference:target-setter-app}}
\noindent\sphinxincludegraphics{{app_architecture-Page-1}.png}


\subsection{r2\_src}
\label{\detokenize{technical_reference:r2-src}}
\noindent\sphinxincludegraphics{{ros2_architecture}.png}


\section{Waypoint queue and state machine}
\label{\detokenize{technical_reference:waypoint-queue-and-state-machine}}

\subsection{Waypoint queue}
\label{\detokenize{technical_reference:waypoint-queue}}\begin{itemize}
\item {} 
\sphinxAtStartPar
\sphinxstylestrong{Waypoint queue:} stores received waypoints from \sphinxtitleref{/waypoint}

\item {} 
\sphinxAtStartPar
\sphinxstylestrong{Active target:} currently executing waypoint

\item {} 
\sphinxAtStartPar
\sphinxstylestrong{Visited waypoints:} stack of completed waypoints for return

\end{itemize}


\section{State machine}
\label{\detokenize{technical_reference:state-machine}}\begin{itemize}
\item {} 
\sphinxAtStartPar
\sphinxcode{\sphinxupquote{IDLE}}: no active target, waiting for waypoint or return command

\item {} 
\sphinxAtStartPar
\sphinxcode{\sphinxupquote{NAVIGATING}}: moving to \sphinxcode{\sphinxupquote{active\_target}}

\item {} 
\sphinxAtStartPar
\sphinxcode{\sphinxupquote{EDIT}}: active target is being edited

\item {} 
\sphinxAtStartPar
\sphinxcode{\sphinxupquote{PAUSED}}: robot is paused by keyboard or automatic pause

\item {} 
\sphinxAtStartPar
\sphinxcode{\sphinxupquote{RETURNING}}: moving to last visited waypoint (return command)

\end{itemize}


\section{Communication protocol}
\label{\detokenize{technical_reference:communication-protocol}}

\subsection{UDP packet structure}
\label{\detokenize{technical_reference:udp-packet-structure}}

\subsubsection{Waypoints packet}
\label{\detokenize{technical_reference:waypoints-packet}}\begin{itemize}
\item {} 
\sphinxAtStartPar
\sphinxstylestrong{Purpose:} send a list of waypoints for the robot to follow

\item {} 
\sphinxAtStartPar
\sphinxstylestrong{Frequency:} send when pressing Send in the app

\end{itemize}

\sphinxAtStartPar
\sphinxstylestrong{Packet Structure (binary):}


\begin{savenotes}\sphinxattablestart
\centering
\begin{tabulary}{\linewidth}[t]{|T|T|T|}
\hline
\sphinxstyletheadfamily 
\sphinxAtStartPar
Offset
&\sphinxstyletheadfamily 
\sphinxAtStartPar
Size (bytes)
&\sphinxstyletheadfamily 
\sphinxAtStartPar
Description
\\
\hline
\sphinxAtStartPar
0
&
\sphinxAtStartPar
4
&
\sphinxAtStartPar
Header (0xAA)
\\
\hline
\sphinxAtStartPar
4
&
\sphinxAtStartPar
4
&
\sphinxAtStartPar
Number of waypoints (counter)
\\
\hline
\sphinxAtStartPar
8
&
\sphinxAtStartPar
4
&
\sphinxAtStartPar
Plan ID (version)
\\
\hline
\sphinxAtStartPar
12
&
\sphinxAtStartPar
16 * N
&
\sphinxAtStartPar
Waypoint (X, Y) in double (8 B)
\\
\hline
\end{tabulary}
\par
\sphinxattableend\end{savenotes}
\begin{itemize}
\item {} 
\sphinxAtStartPar
N = number of waypoints

\item {} 
\sphinxAtStartPar
Coordinates are in meters, converted from percentage position

\item {} 
\sphinxAtStartPar
Byte order: Little Endian

\end{itemize}


\subsubsection{Edit waypoint packet}
\label{\detokenize{technical_reference:edit-waypoint-packet}}\begin{itemize}
\item {} 
\sphinxAtStartPar
\sphinxstylestrong{Purpose:} update a previously sent waypoint

\item {} 
\sphinxAtStartPar
\sphinxstylestrong{Frequency:} send when user edits a waypoint via the app

\end{itemize}

\sphinxAtStartPar
\sphinxstylestrong{Packet structure (binary):}


\begin{savenotes}\sphinxattablestart
\centering
\begin{tabulary}{\linewidth}[t]{|T|T|T|}
\hline
\sphinxstyletheadfamily 
\sphinxAtStartPar
Offset
&\sphinxstyletheadfamily 
\sphinxAtStartPar
Size (bytes)
&\sphinxstyletheadfamily 
\sphinxAtStartPar
Description
\\
\hline
\sphinxAtStartPar
0
&
\sphinxAtStartPar
1
&
\sphinxAtStartPar
Edited flag (1 = edited)
\\
\hline
\sphinxAtStartPar
1
&
\sphinxAtStartPar
4
&
\sphinxAtStartPar
Waypoint index (1\sphinxhyphen{}based)
\\
\hline
\sphinxAtStartPar
5
&
\sphinxAtStartPar
4
&
\sphinxAtStartPar
Plan ID
\\
\hline
\sphinxAtStartPar
9
&
\sphinxAtStartPar
8
&
\sphinxAtStartPar
New X coordinate (double)
\\
\hline
\sphinxAtStartPar
17
&
\sphinxAtStartPar
8
&
\sphinxAtStartPar
New Y coordinate (double)
\\
\hline
\end{tabulary}
\par
\sphinxattableend\end{savenotes}


\subsubsection{Return message packet}
\label{\detokenize{technical_reference:return-message-packet}}\begin{itemize}
\item {} 
\sphinxAtStartPar
\sphinxstylestrong{Purpose:} send a single boolean signal to the robot

\item {} 
\sphinxAtStartPar
\sphinxstylestrong{Frequency:} send when return button is pressed

\end{itemize}

\sphinxAtStartPar
\sphinxstylestrong{Packet structure (binary):}


\begin{savenotes}\sphinxattablestart
\centering
\begin{tabulary}{\linewidth}[t]{|T|T|T|}
\hline
\sphinxstyletheadfamily 
\sphinxAtStartPar
Offset
&\sphinxstyletheadfamily 
\sphinxAtStartPar
Size (bytes)
&\sphinxstyletheadfamily 
\sphinxAtStartPar
Description
\\
\hline
\sphinxAtStartPar
0
&
\sphinxAtStartPar
1
&
\sphinxAtStartPar
Edited flag (0=false,1=true)
\\
\hline
\end{tabulary}
\par
\sphinxattableend\end{savenotes}


\subsubsection{Odometry data reception}
\label{\detokenize{technical_reference:odometry-data-reception}}\begin{itemize}
\item {} 
\sphinxAtStartPar
\sphinxstylestrong{Purpose:} receives current robot odometry

\item {} 
\sphinxAtStartPar
\sphinxstylestrong{Frequency:} 10Hz

\end{itemize}

\sphinxAtStartPar
\sphinxstylestrong{Packet structure (binary):}


\begin{savenotes}\sphinxattablestart
\centering
\begin{tabulary}{\linewidth}[t]{|T|T|T|}
\hline
\sphinxstyletheadfamily 
\sphinxAtStartPar
Offset
&\sphinxstyletheadfamily 
\sphinxAtStartPar
Size (bytes)
&\sphinxstyletheadfamily 
\sphinxAtStartPar
Description
\\
\hline
\sphinxAtStartPar
0
&
\sphinxAtStartPar
8
&
\sphinxAtStartPar
X position (double)
\\
\hline
\sphinxAtStartPar
8
&
\sphinxAtStartPar
8
&
\sphinxAtStartPar
Y position (double)
\\
\hline
\sphinxAtStartPar
16
&
\sphinxAtStartPar
8
&
\sphinxAtStartPar
Yaw (double, radians)
\\
\hline
\end{tabulary}
\par
\sphinxattableend\end{savenotes}
\begin{itemize}
\item {} 
\sphinxAtStartPar
Byte order: Little Endian

\item {} 
\sphinxAtStartPar
Latest odometry stored in memory for waypoint calculations

\end{itemize}


\subsection{UDP data flow}
\label{\detokenize{technical_reference:udp-data-flow}}

\subsubsection{Sending data}
\label{\detokenize{technical_reference:sending-data}}\begin{itemize}
\item {} 
\sphinxAtStartPar
App uses \sphinxcode{\sphinxupquote{NetworkManager.sendWaypoints()}}, \sphinxcode{\sphinxupquote{sendEdit()}}, and \sphinxcode{\sphinxupquote{sendMesage()}}

\item {} 
\sphinxAtStartPar
Waypoints are converted to binary packets before sending

\item {} 
\sphinxAtStartPar
Each send operation is done in a separate thread

\item {} 
\sphinxAtStartPar
Plan ID differentiates waypoint plans

\item {} 
\sphinxAtStartPar
Packets sent via UDP to robot IP/port (default 5050)

\end{itemize}


\subsubsection{Receiving data}
\label{\detokenize{technical_reference:receiving-data}}\begin{itemize}
\item {} 
\sphinxAtStartPar
Persistent UDP socket listens (\sphinxcode{\sphinxupquote{startReceiving(port)}})

\item {} 
\sphinxAtStartPar
Robot sends odometry at 10Hz (X, Y, Yaw)

\item {} 
\sphinxAtStartPar
Incoming packets converted using Little Endian byte order

\item {} 
\sphinxAtStartPar
Latest odometry stored in \sphinxcode{\sphinxupquote{latestOdometry}}

\end{itemize}


\subsubsection{Processing data}
\label{\detokenize{technical_reference:processing-data}}
\sphinxAtStartPar
When adding waypoint:
\begin{itemize}
\item {} 
\sphinxAtStartPar
Convert touch position to normalized coordinates

\item {} 
\sphinxAtStartPar
Compute target X, Y in meters

\item {} 
\sphinxAtStartPar
Add to waypoint list, update UI

\item {} 
\sphinxAtStartPar
Send waypoint via UDP

\end{itemize}

\sphinxAtStartPar
When editing waypoint:
\sphinxhyphen{} Specify index and new coordinates
\sphinxhyphen{} Update list and redraw UI
\sphinxhyphen{} Send edit packet via UDP with plan ID

\sphinxAtStartPar
Return message:
\sphinxhyphen{} Single\sphinxhyphen{}byte packet sent to trigger return behavior

\sphinxAtStartPar
Odometry:
\sphinxhyphen{} Display robot position
\sphinxhyphen{} Compute accurate target positions


\section{Coordinate convention and control logic}
\label{\detokenize{technical_reference:coordinate-convention-and-control-logic}}

\subsection{Coordinate flow}
\label{\detokenize{technical_reference:coordinate-flow}}
\sphinxAtStartPar
Touch \(\rightarrow\) UI \(\rightarrow\) App \(\rightarrow\) Robot

\begin{sphinxVerbatim}[commandchars=\\\{\}]
Touch Event (screen pixels)
   ↓ absolutePosition()
Absolute Coordinates relative to boundaryView
   ↓ normalizedPosition()
Normalized Coordinates (0\textendash{}1)
   ↓ percentagePosition()
Percentage Coordinates (0\textendash{}100\PYGZpc{})
   ↓ targetPosition()
Global Frame (meters) \(\rightarrow\) UDP \(\rightarrow\) Robot
\end{sphinxVerbatim}

\sphinxAtStartPar
Robot \(\rightarrow\) App (Odometry)

\begin{sphinxVerbatim}[commandchars=\\\{\}]
Odometry Packet (X, Y, yaw in meters) \(\rightarrow\) inputPosition()
   ↓ metersToPxX/Y
Pixels on UI \(\rightarrow\) Display robot position
\end{sphinxVerbatim}


\subsection{Frame definition}
\label{\detokenize{technical_reference:frame-definition}}
\sphinxAtStartPar
Local frame:
\begin{itemize}
\item {} 
\sphinxAtStartPar
Origin: the center of the robot

\item {} 
\sphinxAtStartPar
+x: forward

\item {} 
\sphinxAtStartPar
+y: left

\item {} 
\sphinxAtStartPar
+yaw: counter\sphinxhyphen{}clockwise

\end{itemize}

\sphinxAtStartPar
Global frame:
\sphinxhyphen{} Origin: top\sphinxhyphen{}left of the screen on the target\_setter app
\sphinxhyphen{} +x: rightward on the screen on the target\_setter app
\sphinxhyphen{} +y: downward on the screen on the target\_setter app
\sphinxhyphen{} Waypoints that is sent to the robot is in global frame

\sphinxAtStartPar
\sphinxstylestrong{Units:}

\sphinxAtStartPar
On the robot and the data sent from app:
\begin{itemize}
\item {} 
\sphinxAtStartPar
x: meters

\item {} 
\sphinxAtStartPar
y: meters

\item {} 
\sphinxAtStartPar
yaw: radians (data from app only send x and y, yaw must be infer from them)

\item {} 
\sphinxAtStartPar
Vx: m/s

\item {} 
\sphinxAtStartPar
Vy: m/s

\item {} 
\sphinxAtStartPar
Wz: rad/s

\end{itemize}

\sphinxAtStartPar
On the app:
\sphinxhyphen{} x: percentage
\sphinxhyphen{} y: percentage


\subsection{Touch input conversion}
\label{\detokenize{technical_reference:touch-input-conversion}}
\sphinxAtStartPar
\sphinxstylestrong{Absolute position (touch):}

\begin{sphinxVerbatim}[commandchars=\\\{\}]
\PYG{k+kt}{double}\PYG{+w}{ }\PYG{n}{absoluteX}\PYG{+w}{ }\PYG{o}{=}\PYG{+w}{ }\PYG{n}{rawX}\PYG{+w}{ }\PYG{o}{\PYGZhy{}}\PYG{+w}{ }\PYG{n}{boundaryLeft}\PYG{p}{;}
\PYG{k+kt}{double}\PYG{+w}{ }\PYG{n}{absoluteY}\PYG{+w}{ }\PYG{o}{=}\PYG{+w}{ }\PYG{n}{rawY}\PYG{+w}{ }\PYG{o}{\PYGZhy{}}\PYG{+w}{ }\PYG{n}{boundaryTop}\PYG{p}{;}
\end{sphinxVerbatim}
\begin{itemize}
\item {} 
\sphinxAtStartPar
Converts raw screen coordinate to coordinate relative to \sphinxtitleref{boundaryView}

\item {} 
\sphinxAtStartPar
Clamped to boundary edges

\end{itemize}

\sphinxAtStartPar
\sphinxstylestrong{Normalized position (0\sphinxhyphen{}1):}

\begin{sphinxVerbatim}[commandchars=\\\{\}]
\PYG{k+kt}{double}\PYG{+w}{ }\PYG{n}{normalizedX}\PYG{+w}{ }\PYG{o}{=}\PYG{+w}{ }\PYG{n}{absoluteX}\PYG{+w}{ }\PYG{o}{/}\PYG{+w}{ }\PYG{n}{screenWidth}\PYG{p}{;}
\PYG{k+kt}{double}\PYG{+w}{ }\PYG{n}{normalizedY}\PYG{+w}{ }\PYG{o}{=}\PYG{+w}{ }\PYG{n}{absoluteY}\PYG{+w}{ }\PYG{o}{/}\PYG{+w}{ }\PYG{n}{screenHeight}\PYG{p}{;}
\end{sphinxVerbatim}
\begin{itemize}
\item {} 
\sphinxAtStartPar
Represents position as fraction of the screen

\end{itemize}

\sphinxAtStartPar
\sphinxstylestrong{Percentage position (0\sphinxhyphen{}100\%):}

\begin{sphinxVerbatim}[commandchars=\\\{\}]
\PYG{k+kt}{double}\PYG{+w}{ }\PYG{n}{percentageX}\PYG{+w}{ }\PYG{o}{=}\PYG{+w}{ }\PYG{n}{normalizedX}\PYG{+w}{ }\PYG{o}{*}\PYG{+w}{ }\PYG{l+m+mf}{100.0}\PYG{p}{;}
\PYG{k+kt}{double}\PYG{+w}{ }\PYG{n}{percentageY}\PYG{+w}{ }\PYG{o}{=}\PYG{+w}{ }\PYG{n}{normalizedY}\PYG{+w}{ }\PYG{o}{*}\PYG{+w}{ }\PYG{l+m+mf}{100.0}\PYG{p}{;}
\end{sphinxVerbatim}
\begin{itemize}
\item {} 
\sphinxAtStartPar
Used as intermediate format for calculation into world meters

\item {} 
\sphinxAtStartPar
Checked against 0\sphinxhyphen{}100\% boundary

\end{itemize}

\sphinxAtStartPar
\sphinxstylestrong{Global meters (for UDP):}

\begin{sphinxVerbatim}[commandchars=\\\{\}]
\PYG{k+kt}{double}\PYG{+w}{ }\PYG{n}{targetX}\PYG{+w}{ }\PYG{o}{=}\PYG{+w}{ }\PYG{p}{(}\PYG{n}{percentageX}\PYG{+w}{ }\PYG{o}{/}\PYG{+w}{ }\PYG{l+m+mf}{100.0}\PYG{p}{)}\PYG{+w}{ }\PYG{o}{*}\PYG{+w}{ }\PYG{n}{GAME\PYGZus{}FIELD\PYGZus{}X}\PYG{p}{;}
\PYG{k+kt}{double}\PYG{+w}{ }\PYG{n}{targetY}\PYG{+w}{ }\PYG{o}{=}\PYG{+w}{ }\PYG{p}{(}\PYG{n}{percentageY}\PYG{+w}{ }\PYG{o}{/}\PYG{+w}{ }\PYG{l+m+mf}{100.0}\PYG{p}{)}\PYG{+w}{ }\PYG{o}{*}\PYG{+w}{ }\PYG{n}{GAME\PYGZus{}FIELD\PYGZus{}Y}\PYG{p}{;}
\end{sphinxVerbatim}
\begin{itemize}
\item {} 
\sphinxAtStartPar
Converts to meters based on field dimension

\item {} 
\sphinxAtStartPar
Result is sent to robot

\end{itemize}

\sphinxAtStartPar
\sphinxstylestrong{Robot odometry (UI display):}

\begin{sphinxVerbatim}[commandchars=\\\{\}]
\PYG{k+kt}{double}\PYG{+w}{ }\PYG{n}{inputX}\PYG{+w}{ }\PYG{o}{=}\PYG{+w}{ }\PYG{n}{odometryData}\PYG{o}{[}\PYG{l+m+mi}{0}\PYG{o}{]}\PYG{+w}{ }\PYG{o}{*}\PYG{+w}{ }\PYG{p}{(}\PYG{n}{screenWidth}\PYG{+w}{ }\PYG{o}{/}\PYG{+w}{ }\PYG{n}{GAME\PYGZus{}FIELD\PYGZus{}X}\PYG{p}{)}\PYG{p}{;}
\PYG{k+kt}{double}\PYG{+w}{ }\PYG{n}{inputY}\PYG{+w}{ }\PYG{o}{=}\PYG{+w}{ }\PYG{n}{odometryData}\PYG{o}{[}\PYG{l+m+mi}{1}\PYG{o}{]}\PYG{+w}{ }\PYG{o}{*}\PYG{+w}{ }\PYG{p}{(}\PYG{n}{screenHeight}\PYG{+w}{ }\PYG{o}{/}\PYG{+w}{ }\PYG{n}{GAME\PYGZus{}FIELD\PYGZus{}Y}\PYG{p}{)}\PYG{p}{;}
\end{sphinxVerbatim}
\begin{itemize}
\item {} 
\sphinxAtStartPar
Converts meters to pixels for UI display

\item {} 
\sphinxAtStartPar
Clamped boundary using \sphinxtitleref{validateBoundary()}

\end{itemize}


\section{Control logic}
\label{\detokenize{technical_reference:control-logic}}

\subsection{PD controller}
\label{\detokenize{technical_reference:pd-controller}}\begin{itemize}
\item {} 
\sphinxAtStartPar
\sphinxstylestrong{Position errors:}

\end{itemize}

\sphinxAtStartPar
computed in global frame

\begin{sphinxVerbatim}[commandchars=\\\{\}]
\PYG{n}{current\PYGZus{}x\PYGZus{}error\PYGZus{}g} \PYG{o}{=} \PYG{n+nb+bp}{self}\PYG{o}{.}\PYG{n}{desired\PYGZus{}x} \PYG{o}{\PYGZhy{}} \PYG{n+nb+bp}{self}\PYG{o}{.}\PYG{n}{current\PYGZus{}x}
\PYG{n}{current\PYGZus{}y\PYGZus{}error\PYGZus{}g} \PYG{o}{=} \PYG{n+nb+bp}{self}\PYG{o}{.}\PYG{n}{desired\PYGZus{}y} \PYG{o}{\PYGZhy{}} \PYG{n+nb+bp}{self}\PYG{o}{.}\PYG{n}{current\PYGZus{}y}

\PYG{n}{current\PYGZus{}yaw\PYGZus{}p\PYGZus{}error} \PYG{o}{=} \PYG{n+nb+bp}{self}\PYG{o}{.}\PYG{n}{normalize\PYGZus{}angle}\PYG{p}{(}\PYG{n+nb+bp}{self}\PYG{o}{.}\PYG{n}{desired\PYGZus{}yaw} \PYG{o}{\PYGZhy{}} \PYG{n+nb+bp}{self}\PYG{o}{.}\PYG{n}{current\PYGZus{}yaw}\PYG{p}{)}
\end{sphinxVerbatim}
\begin{itemize}
\item {} 
\sphinxAtStartPar
\sphinxstylestrong{Derivative errors:}

\end{itemize}

\begin{sphinxVerbatim}[commandchars=\\\{\}]
\PYG{n}{current\PYGZus{}yaw\PYGZus{}d\PYGZus{}error} \PYG{o}{=} \PYG{p}{(}\PYG{n}{current\PYGZus{}yaw\PYGZus{}p\PYGZus{}error} \PYG{o}{\PYGZhy{}} \PYG{n+nb+bp}{self}\PYG{o}{.}\PYG{n}{previous\PYGZus{}yaw\PYGZus{}p\PYGZus{}error}\PYG{p}{)}\PYG{o}{/}\PYG{n+nb+bp}{self}\PYG{o}{.}\PYG{n}{dt}
\PYG{n}{current\PYGZus{}x\PYGZus{}d\PYGZus{}error} \PYG{o}{=} \PYG{p}{(}\PYG{n}{current\PYGZus{}x\PYGZus{}error\PYGZus{}g} \PYG{o}{\PYGZhy{}} \PYG{n+nb+bp}{self}\PYG{o}{.}\PYG{n}{previous\PYGZus{}x\PYGZus{}p\PYGZus{}error}\PYG{p}{)}\PYG{o}{/}\PYG{n+nb+bp}{self}\PYG{o}{.}\PYG{n}{dt}
\PYG{n}{current\PYGZus{}y\PYGZus{}d\PYGZus{}error} \PYG{o}{=} \PYG{p}{(}\PYG{n}{current\PYGZus{}y\PYGZus{}error\PYGZus{}g} \PYG{o}{\PYGZhy{}} \PYG{n+nb+bp}{self}\PYG{o}{.}\PYG{n}{previous\PYGZus{}y\PYGZus{}p\PYGZus{}error}\PYG{p}{)}\PYG{o}{/}\PYG{n+nb+bp}{self}\PYG{o}{.}\PYG{n}{dt}
\end{sphinxVerbatim}
\begin{itemize}
\item {} 
\sphinxAtStartPar
\sphinxstylestrong{Complementary filtering:}

\end{itemize}

\sphinxAtStartPar
apply complementary filter on D term to smooth the movement

\begin{sphinxVerbatim}[commandchars=\\\{\}]
\PYG{n}{current\PYGZus{}x\PYGZus{}d\PYGZus{}error} \PYG{o}{=} \PYG{n}{ALPHA\PYGZus{}D} \PYG{o}{*} \PYG{n}{current\PYGZus{}x\PYGZus{}d\PYGZus{}error} \PYG{o}{+} \PYG{p}{(}\PYG{l+m+mi}{1} \PYG{o}{\PYGZhy{}} \PYG{n}{ALPHA\PYGZus{}D}\PYG{p}{)} \PYG{o}{*} \PYG{n+nb+bp}{self}\PYG{o}{.}\PYG{n}{previous\PYGZus{}x\PYGZus{}d\PYGZus{}error}
\PYG{n}{current\PYGZus{}y\PYGZus{}d\PYGZus{}error} \PYG{o}{=} \PYG{n}{ALPHA\PYGZus{}D} \PYG{o}{*} \PYG{n}{current\PYGZus{}y\PYGZus{}d\PYGZus{}error} \PYG{o}{+} \PYG{p}{(}\PYG{l+m+mi}{1} \PYG{o}{\PYGZhy{}} \PYG{n}{ALPHA\PYGZus{}D}\PYG{p}{)} \PYG{o}{*} \PYG{n+nb+bp}{self}\PYG{o}{.}\PYG{n}{previous\PYGZus{}y\PYGZus{}d\PYGZus{}error}
\PYG{n}{current\PYGZus{}yaw\PYGZus{}d\PYGZus{}error} \PYG{o}{=} \PYG{n}{ALPHA\PYGZus{}D} \PYG{o}{*} \PYG{n}{current\PYGZus{}yaw\PYGZus{}d\PYGZus{}error} \PYG{o}{+} \PYG{p}{(}\PYG{l+m+mi}{1} \PYG{o}{\PYGZhy{}} \PYG{n}{ALPHA\PYGZus{}D}\PYG{p}{)} \PYG{o}{*} \PYG{n+nb+bp}{self}\PYG{o}{.}\PYG{n}{previous\PYGZus{}yaw\PYGZus{}d\PYGZus{}error}
\end{sphinxVerbatim}
\begin{itemize}
\item {} 
\sphinxAtStartPar
\sphinxstylestrong{PD controller:}

\end{itemize}

\begin{sphinxVerbatim}[commandchars=\\\{\}]
\PYG{n}{linear\PYGZus{}vel\PYGZus{}x\PYGZus{}g} \PYG{o}{=} \PYG{n+nb+bp}{self}\PYG{o}{.}\PYG{n}{k\PYGZus{}p\PYGZus{}linear} \PYG{o}{*} \PYG{n}{current\PYGZus{}x\PYGZus{}error\PYGZus{}g} \PYG{o}{+} \PYG{n+nb+bp}{self}\PYG{o}{.}\PYG{n}{k\PYGZus{}d} \PYG{o}{*} \PYG{n}{current\PYGZus{}x\PYGZus{}d\PYGZus{}error}
\PYG{n}{linear\PYGZus{}vel\PYGZus{}y\PYGZus{}g} \PYG{o}{=} \PYG{n+nb+bp}{self}\PYG{o}{.}\PYG{n}{k\PYGZus{}p\PYGZus{}linear} \PYG{o}{*} \PYG{n}{current\PYGZus{}y\PYGZus{}error\PYGZus{}g} \PYG{o}{+} \PYG{n+nb+bp}{self}\PYG{o}{.}\PYG{n}{k\PYGZus{}d} \PYG{o}{*} \PYG{n}{current\PYGZus{}y\PYGZus{}d\PYGZus{}error}
\PYG{n}{angular\PYGZus{}vel\PYGZus{}z\PYGZus{}g} \PYG{o}{=} \PYG{n+nb+bp}{self}\PYG{o}{.}\PYG{n}{k\PYGZus{}p\PYGZus{}yaw} \PYG{o}{*} \PYG{n}{current\PYGZus{}yaw\PYGZus{}p\PYGZus{}error} \PYG{o}{+} \PYG{n+nb+bp}{self}\PYG{o}{.}\PYG{n}{k\PYGZus{}d\PYGZus{}yaw} \PYG{o}{*} \PYG{n}{current\PYGZus{}yaw\PYGZus{}d\PYGZus{}error}
\end{sphinxVerbatim}
\begin{itemize}
\item {} 
\sphinxAtStartPar
\sphinxstylestrong{Frame conversion:}

\end{itemize}

\begin{sphinxVerbatim}[commandchars=\\\{\}]
\PYG{n}{linear\PYGZus{}vel\PYGZus{}x\PYGZus{}l} \PYG{o}{=}  \PYG{n}{math}\PYG{o}{.}\PYG{n}{cos}\PYG{p}{(}\PYG{n+nb+bp}{self}\PYG{o}{.}\PYG{n}{current\PYGZus{}yaw}\PYG{p}{)} \PYG{o}{*} \PYG{n}{current\PYGZus{}x\PYGZus{}error\PYGZus{}g} \PYG{o}{+} \PYG{n}{math}\PYG{o}{.}\PYG{n}{sin}\PYG{p}{(}\PYG{n+nb+bp}{self}\PYG{o}{.}\PYG{n}{current\PYGZus{}yaw}\PYG{p}{)} \PYG{o}{*} \PYG{n}{current\PYGZus{}y\PYGZus{}error\PYGZus{}g}
\PYG{n}{linear\PYGZus{}vel\PYGZus{}y\PYGZus{}l} \PYG{o}{=} \PYG{o}{\PYGZhy{}}\PYG{n}{math}\PYG{o}{.}\PYG{n}{sin}\PYG{p}{(}\PYG{n+nb+bp}{self}\PYG{o}{.}\PYG{n}{current\PYGZus{}yaw}\PYG{p}{)} \PYG{o}{*} \PYG{n}{current\PYGZus{}x\PYGZus{}error\PYGZus{}g} \PYG{o}{+} \PYG{n}{math}\PYG{o}{.}\PYG{n}{cos}\PYG{p}{(}\PYG{n+nb+bp}{self}\PYG{o}{.}\PYG{n}{current\PYGZus{}yaw}\PYG{p}{)} \PYG{o}{*} \PYG{n}{current\PYGZus{}y\PYGZus{}error\PYGZus{}g}
\end{sphinxVerbatim}
\begin{itemize}
\item {} 
\sphinxAtStartPar
Units must match: meters (x, y), radians (yaw), m/s (velocities)

\item {} 
\sphinxAtStartPar
Linear velocity in robot frame, angular velocity around z\sphinxhyphen{}axis

\item {} 
\sphinxAtStartPar
Low\sphinxhyphen{}pass filtering applied on derivative term

\item {} 
\sphinxAtStartPar
Return \sphinxtitleref{linear\_vel\_x}, \sphinxtitleref{linear\_vel\_y}, and \sphinxtitleref{angular\_vel\_z}

\end{itemize}


\subsection{Dead reckoning}
\label{\detokenize{technical_reference:dead-reckoning}}
\sphinxAtStartPar
If the velocity is less than error threshold, overwrite it with 0


\subsection{Speed limit}
\label{\detokenize{technical_reference:speed-limit}}
\sphinxAtStartPar
If velocity exceeds maximum speed, overwrite it with maximum speed


\section{Path planning}
\label{\detokenize{technical_reference:path-planning}}

\subsection{Arrivel detection}
\label{\detokenize{technical_reference:arrivel-detection}}\begin{itemize}
\item {} 
\sphinxAtStartPar
Target reached if:
\sphinxhyphen{} Distance \textless{} \sphinxcode{\sphinxupquote{goal\_tolerance\_pos = 0.05m}}
\sphinxhyphen{} Remains withitn tolerance for \sphinxcode{\sphinxupquote{0.5s}}

\item {} 
\sphinxAtStartPar
Arrival time accumulated using \sphinxcode{\sphinxupquote{dt}} from timer callback

\end{itemize}


\subsection{Waypoint callback}
\label{\detokenize{technical_reference:waypoint-callback}}\begin{itemize}
\item {} 
\sphinxAtStartPar
Store each waypoint in a dictionary

\item {} 
\sphinxAtStartPar
Append it into the queue

\item {} 
\sphinxAtStartPar
Copy the waypoint in the queue to \sphinxcode{\sphinxupquote{active\_target}} to execute

\end{itemize}


\subsection{Waypoints execution}
\label{\detokenize{technical_reference:waypoints-execution}}\begin{itemize}
\item {} 
\sphinxAtStartPar
if \sphinxcode{\sphinxupquote{active\_target}} does not exist, and there’s waypoint in the queue, pop the queue, then reset robot state to go to the waypoint that’s being popped

\item {} 
\sphinxAtStartPar
if return is triggered, and there’s visited waypoint, command the robot to go the visited waypoint

\item {} 
\sphinxAtStartPar
Otherwise, stop the robot

\end{itemize}


\subsection{Waypoint update}
\label{\detokenize{technical_reference:waypoint-update}}\begin{itemize}
\item {} 
\sphinxAtStartPar
Store the edited waypoint in a dictionary

\item {} 
\sphinxAtStartPar
If the edited waypoint is active, use complementary to smooth the robot movement when update to prevent jumping suddenly

\end{itemize}

\begin{sphinxVerbatim}[commandchars=\\\{\}]
\PYG{n+nb+bp}{self}\PYG{o}{.}\PYG{n}{active\PYGZus{}target}\PYG{p}{[}\PYG{l+s+s1}{\PYGZsq{}}\PYG{l+s+s1}{x}\PYG{l+s+s1}{\PYGZsq{}}\PYG{p}{]} \PYG{o}{=} \PYG{p}{(}\PYG{l+m+mi}{1} \PYG{o}{\PYGZhy{}} \PYG{n}{ALPHA}\PYG{p}{)} \PYG{o}{*} \PYG{n+nb+bp}{self}\PYG{o}{.}\PYG{n}{active\PYGZus{}target}\PYG{p}{[}\PYG{l+s+s1}{\PYGZsq{}}\PYG{l+s+s1}{x}\PYG{l+s+s1}{\PYGZsq{}}\PYG{p}{]} \PYG{o}{+} \PYG{n}{ALPHA} \PYG{o}{*} \PYG{n}{update\PYGZus{}wp}\PYG{p}{[}\PYG{l+s+s1}{\PYGZsq{}}\PYG{l+s+s1}{updated\PYGZus{}x}\PYG{l+s+s1}{\PYGZsq{}}\PYG{p}{]}
\PYG{n+nb+bp}{self}\PYG{o}{.}\PYG{n}{active\PYGZus{}target}\PYG{p}{[}\PYG{l+s+s1}{\PYGZsq{}}\PYG{l+s+s1}{y}\PYG{l+s+s1}{\PYGZsq{}}\PYG{p}{]} \PYG{o}{=} \PYG{p}{(}\PYG{l+m+mi}{1} \PYG{o}{\PYGZhy{}} \PYG{n}{ALPHA}\PYG{p}{)} \PYG{o}{*} \PYG{n+nb+bp}{self}\PYG{o}{.}\PYG{n}{active\PYGZus{}target}\PYG{p}{[}\PYG{l+s+s1}{\PYGZsq{}}\PYG{l+s+s1}{y}\PYG{l+s+s1}{\PYGZsq{}}\PYG{p}{]} \PYG{o}{+} \PYG{n}{ALPHA} \PYG{o}{*} \PYG{n}{update\PYGZus{}wp}\PYG{p}{[}\PYG{l+s+s1}{\PYGZsq{}}\PYG{l+s+s1}{updated\PYGZus{}y}\PYG{l+s+s1}{\PYGZsq{}}\PYG{p}{]}
\end{sphinxVerbatim}
\begin{itemize}
\item {} 
\sphinxAtStartPar
Otherwise, just update the waypoint as normal

\end{itemize}


\subsection{Pause}
\label{\detokenize{technical_reference:pause}}\begin{itemize}
\item {} 
\sphinxAtStartPar
Robot pauses motion at each target

\item {} 
\sphinxAtStartPar
If distance from the robot to the target is less than tolerance, stop the robot

\item {} 
\sphinxAtStartPar
If the elapsed time during pause exceeds pause duration, command the robot to go to the next waypoint

\end{itemize}


\subsection{Return}
\label{\detokenize{technical_reference:return}}\begin{itemize}
\item {} 
\sphinxAtStartPar
Update the return flag

\end{itemize}


\section{Odometry handling}
\label{\detokenize{technical_reference:odometry-handling}}

\subsection{Frame conversion}
\label{\detokenize{technical_reference:frame-conversion}}\begin{itemize}
\item {} 
\sphinxAtStartPar
Calculate velocity in local frame

\end{itemize}

\begin{sphinxVerbatim}[commandchars=\\\{\}]
\PYG{n}{linear\PYGZus{}vel\PYGZus{}x\PYGZus{}local} \PYG{o}{=} \PYG{p}{(}\PYG{n+nb+bp}{self}\PYG{o}{.}\PYG{n}{R}\PYG{o}{/}\PYG{l+m+mi}{4}\PYG{p}{)} \PYG{o}{*} \PYG{p}{(}\PYG{n+nb+bp}{self}\PYG{o}{.}\PYG{n}{motor\PYGZus{}vel}\PYG{p}{[}\PYG{l+m+mi}{0}\PYG{p}{]} \PYG{o}{+} \PYG{n+nb+bp}{self}\PYG{o}{.}\PYG{n}{motor\PYGZus{}vel}\PYG{p}{[}\PYG{l+m+mi}{1}\PYG{p}{]} \PYG{o}{+} \PYG{n+nb+bp}{self}\PYG{o}{.}\PYG{n}{motor\PYGZus{}vel}\PYG{p}{[}\PYG{l+m+mi}{2}\PYG{p}{]} \PYG{o}{+} \PYG{n+nb+bp}{self}\PYG{o}{.}\PYG{n}{motor\PYGZus{}vel}\PYG{p}{[}\PYG{l+m+mi}{3}\PYG{p}{]}\PYG{p}{)}
\PYG{n}{linear\PYGZus{}vel\PYGZus{}y\PYGZus{}local} \PYG{o}{=} \PYG{p}{(}\PYG{n+nb+bp}{self}\PYG{o}{.}\PYG{n}{R}\PYG{o}{/}\PYG{l+m+mi}{4}\PYG{p}{)} \PYG{o}{*} \PYG{p}{(}\PYG{o}{\PYGZhy{}}\PYG{n+nb+bp}{self}\PYG{o}{.}\PYG{n}{motor\PYGZus{}vel}\PYG{p}{[}\PYG{l+m+mi}{0}\PYG{p}{]} \PYG{o}{+} \PYG{n+nb+bp}{self}\PYG{o}{.}\PYG{n}{motor\PYGZus{}vel}\PYG{p}{[}\PYG{l+m+mi}{1}\PYG{p}{]} \PYG{o}{+} \PYG{n+nb+bp}{self}\PYG{o}{.}\PYG{n}{motor\PYGZus{}vel}\PYG{p}{[}\PYG{l+m+mi}{2}\PYG{p}{]} \PYG{o}{\PYGZhy{}} \PYG{n+nb+bp}{self}\PYG{o}{.}\PYG{n}{motor\PYGZus{}vel}\PYG{p}{[}\PYG{l+m+mi}{3}\PYG{p}{]}\PYG{p}{)}
\PYG{n}{w\PYGZus{}z} \PYG{o}{=} \PYG{p}{(}\PYG{n+nb+bp}{self}\PYG{o}{.}\PYG{n}{R} \PYG{o}{*} \PYG{p}{(}\PYG{o}{\PYGZhy{}}\PYG{n+nb+bp}{self}\PYG{o}{.}\PYG{n}{motor\PYGZus{}vel}\PYG{p}{[}\PYG{l+m+mi}{0}\PYG{p}{]} \PYG{o}{+} \PYG{n+nb+bp}{self}\PYG{o}{.}\PYG{n}{motor\PYGZus{}vel}\PYG{p}{[}\PYG{l+m+mi}{1}\PYG{p}{]} \PYG{o}{\PYGZhy{}} \PYG{n+nb+bp}{self}\PYG{o}{.}\PYG{n}{motor\PYGZus{}vel}\PYG{p}{[}\PYG{l+m+mi}{2}\PYG{p}{]} \PYG{o}{+} \PYG{n+nb+bp}{self}\PYG{o}{.}\PYG{n}{motor\PYGZus{}vel}\PYG{p}{[}\PYG{l+m+mi}{3}\PYG{p}{]}\PYG{p}{)}\PYG{p}{)}\PYG{o}{/}\PYG{p}{(}\PYG{l+m+mi}{2} \PYG{o}{*} \PYG{p}{(}\PYG{n+nb+bp}{self}\PYG{o}{.}\PYG{n}{l\PYGZus{}x} \PYG{o}{+} \PYG{n+nb+bp}{self}\PYG{o}{.}\PYG{n}{l\PYGZus{}y}\PYG{p}{)}\PYG{p}{)}
\end{sphinxVerbatim}
\begin{itemize}
\item {} 
\sphinxAtStartPar
Convert them to global frame

\end{itemize}

\begin{sphinxVerbatim}[commandchars=\\\{\}]
\PYG{n}{v\PYGZus{}x\PYGZus{}global} \PYG{o}{=} \PYG{n}{math}\PYG{o}{.}\PYG{n}{cos}\PYG{p}{(}\PYG{n+nb+bp}{self}\PYG{o}{.}\PYG{n}{current\PYGZus{}yaw}\PYG{p}{)} \PYG{o}{*} \PYG{n}{linear\PYGZus{}vel\PYGZus{}x\PYGZus{}local} \PYG{o}{\PYGZhy{}} \PYG{n}{math}\PYG{o}{.}\PYG{n}{sin}\PYG{p}{(}\PYG{n+nb+bp}{self}\PYG{o}{.}\PYG{n}{current\PYGZus{}yaw}\PYG{p}{)} \PYG{o}{*} \PYG{n}{linear\PYGZus{}vel\PYGZus{}y\PYGZus{}local}
\PYG{n}{v\PYGZus{}y\PYGZus{}global} \PYG{o}{=} \PYG{n}{math}\PYG{o}{.}\PYG{n}{sin}\PYG{p}{(}\PYG{n+nb+bp}{self}\PYG{o}{.}\PYG{n}{current\PYGZus{}yaw}\PYG{p}{)} \PYG{o}{*} \PYG{n}{linear\PYGZus{}vel\PYGZus{}x\PYGZus{}local} \PYG{o}{+} \PYG{n}{math}\PYG{o}{.}\PYG{n}{cos}\PYG{p}{(}\PYG{n+nb+bp}{self}\PYG{o}{.}\PYG{n}{current\PYGZus{}yaw}\PYG{p}{)} \PYG{o}{*} \PYG{n}{linear\PYGZus{}vel\PYGZus{}y\PYGZus{}local}
\end{sphinxVerbatim}
\begin{itemize}
\item {} 
\sphinxAtStartPar
Integrate velocity to get position

\end{itemize}

\begin{sphinxVerbatim}[commandchars=\\\{\}]
\PYG{n+nb+bp}{self}\PYG{o}{.}\PYG{n}{x} \PYG{o}{+}\PYG{o}{=} \PYG{n}{v\PYGZus{}x\PYGZus{}global} \PYG{o}{*} \PYG{n+nb+bp}{self}\PYG{o}{.}\PYG{n}{dt}
\PYG{n+nb+bp}{self}\PYG{o}{.}\PYG{n}{y} \PYG{o}{+}\PYG{o}{=} \PYG{n}{v\PYGZus{}y\PYGZus{}global} \PYG{o}{*} \PYG{n+nb+bp}{self}\PYG{o}{.}\PYG{n}{dt}
\end{sphinxVerbatim}


\subsection{Offset}
\label{\detokenize{technical_reference:offset}}
\begin{sphinxVerbatim}[commandchars=\\\{\}]
\PYG{n+nb+bp}{self}\PYG{o}{.}\PYG{n}{x\PYGZus{}odom} \PYG{o}{=} \PYG{p}{(}\PYG{n+nb+bp}{self}\PYG{o}{.}\PYG{n}{x} \PYG{o}{\PYGZhy{}} \PYG{n+nb+bp}{self}\PYG{o}{.}\PYG{n}{x\PYGZus{}start}\PYG{p}{)} \PYG{o}{\PYGZhy{}} \PYG{n+nb+bp}{self}\PYG{o}{.}\PYG{n}{frame\PYGZus{}offset\PYGZus{}x}
\PYG{n+nb+bp}{self}\PYG{o}{.}\PYG{n}{y\PYGZus{}odom} \PYG{o}{=} \PYG{p}{(}\PYG{n+nb+bp}{self}\PYG{o}{.}\PYG{n}{y} \PYG{o}{\PYGZhy{}} \PYG{n+nb+bp}{self}\PYG{o}{.}\PYG{n}{y\PYGZus{}start}\PYG{p}{)} \PYG{o}{\PYGZhy{}} \PYG{n+nb+bp}{self}\PYG{o}{.}\PYG{n}{frame\PYGZus{}offset\PYGZus{}y}
\end{sphinxVerbatim}


\subsection{Yaw}
\label{\detokenize{technical_reference:yaw}}\begin{itemize}
\item {} 
\sphinxAtStartPar
Calculate yaw from quaternion

\end{itemize}

\begin{sphinxVerbatim}[commandchars=\\\{\}]
\PYG{k}{def}\PYG{+w}{ }\PYG{n+nf}{calculate\PYGZus{}yaw}\PYG{p}{(}\PYG{n+nb+bp}{self}\PYG{p}{,} \PYG{n}{q\PYGZus{}x}\PYG{p}{,} \PYG{n}{q\PYGZus{}y}\PYG{p}{,} \PYG{n}{q\PYGZus{}z}\PYG{p}{,} \PYG{n}{q\PYGZus{}w}\PYG{p}{)}\PYG{p}{:}
   \PYG{n}{siny\PYGZus{}cosp} \PYG{o}{=} \PYG{l+m+mi}{2} \PYG{o}{*} \PYG{p}{(}\PYG{n}{q\PYGZus{}w} \PYG{o}{*} \PYG{n}{q\PYGZus{}z} \PYG{o}{+} \PYG{n}{q\PYGZus{}x} \PYG{o}{*} \PYG{n}{q\PYGZus{}y}\PYG{p}{)}
   \PYG{n}{cosy\PYGZus{}sinp} \PYG{o}{=} \PYG{l+m+mi}{1} \PYG{o}{\PYGZhy{}} \PYG{l+m+mi}{2} \PYG{o}{*} \PYG{p}{(}\PYG{n}{q\PYGZus{}y}\PYG{o}{*}\PYG{o}{*}\PYG{l+m+mi}{2} \PYG{o}{+} \PYG{n}{q\PYGZus{}z}\PYG{o}{*}\PYG{o}{*}\PYG{l+m+mi}{2}\PYG{p}{)}
   \PYG{k}{return} \PYG{n}{math}\PYG{o}{.}\PYG{n}{atan2}\PYG{p}{(}\PYG{n}{siny\PYGZus{}cosp}\PYG{p}{,} \PYG{n}{cosy\PYGZus{}sinp}\PYG{p}{)}
\end{sphinxVerbatim}
\begin{itemize}
\item {} 
\sphinxAtStartPar
Normalize yaw

\end{itemize}

\begin{sphinxVerbatim}[commandchars=\\\{\}]
\PYG{k}{def}\PYG{+w}{ }\PYG{n+nf}{normalize\PYGZus{}angle}\PYG{p}{(}\PYG{n+nb+bp}{self}\PYG{p}{,} \PYG{n}{a}\PYG{p}{)}\PYG{p}{:}
   \PYG{k}{return} \PYG{n}{math}\PYG{o}{.}\PYG{n}{atan2}\PYG{p}{(}\PYG{n}{math}\PYG{o}{.}\PYG{n}{sin}\PYG{p}{(}\PYG{n}{a}\PYG{p}{)}\PYG{p}{,} \PYG{n}{math}\PYG{o}{.}\PYG{n}{cos}\PYG{p}{(}\PYG{n}{a}\PYG{p}{)}\PYG{p}{)}
\end{sphinxVerbatim}
\begin{itemize}
\item {} 
\sphinxAtStartPar
Offset yaw

\end{itemize}

\begin{sphinxVerbatim}[commandchars=\\\{\}]
\PYG{n+nb+bp}{self}\PYG{o}{.}\PYG{n}{current\PYGZus{}yaw} \PYG{o}{=} \PYG{p}{(}\PYG{n+nb+bp}{self}\PYG{o}{.}\PYG{n}{yaw} \PYG{o}{\PYGZhy{}} \PYG{n+nb+bp}{self}\PYG{o}{.}\PYG{n}{yaw\PYGZus{}start}\PYG{p}{)} \PYG{o}{+} \PYG{l+m+mi}{2} \PYG{o}{*} \PYG{n}{math}\PYG{o}{.}\PYG{n}{pi} \PYG{o}{*} \PYG{n+nb+bp}{self}\PYG{o}{.}\PYG{n}{overflow\PYGZus{}counter}
\end{sphinxVerbatim}
\begin{itemize}
\item {} 
\sphinxAtStartPar
Convert to quaternion to publish

\end{itemize}

\begin{sphinxVerbatim}[commandchars=\\\{\}]
\PYG{k}{def}\PYG{+w}{ }\PYG{n+nf}{publish\PYGZus{}yaw}\PYG{p}{(}\PYG{n+nb+bp}{self}\PYG{p}{,} \PYG{n}{yaw}\PYG{p}{)}\PYG{p}{:}
   \PYG{n}{q} \PYG{o}{=} \PYG{n}{Quaternion}\PYG{p}{(}\PYG{p}{)}
   \PYG{n}{q}\PYG{o}{.}\PYG{n}{w} \PYG{o}{=} \PYG{n}{math}\PYG{o}{.}\PYG{n}{cos}\PYG{p}{(}\PYG{n}{yaw}\PYG{o}{/}\PYG{l+m+mi}{2}\PYG{p}{)}
   \PYG{n}{q}\PYG{o}{.}\PYG{n}{x} \PYG{o}{=} \PYG{l+m+mf}{0.0}
   \PYG{n}{q}\PYG{o}{.}\PYG{n}{y} \PYG{o}{=} \PYG{l+m+mf}{0.0}
   \PYG{n}{q}\PYG{o}{.}\PYG{n}{z} \PYG{o}{=} \PYG{n}{math}\PYG{o}{.}\PYG{n}{sin}\PYG{p}{(}\PYG{n}{yaw}\PYG{o}{/}\PYG{l+m+mi}{2}\PYG{p}{)}
   \PYG{k}{return} \PYG{n}{q}
\end{sphinxVerbatim}


\section{Inverse Kinematic}
\label{\detokenize{technical_reference:inverse-kinematic}}
\begin{sphinxVerbatim}[commandchars=\\\{\}]
\PYG{n}{V\PYGZus{}wheel} \PYG{o}{=} \PYG{p}{[}
\PYG{n}{Vx}\PYG{o}{/}\PYG{n+nb+bp}{self}\PYG{o}{.}\PYG{n}{R} \PYG{o}{\PYGZhy{}} \PYG{n}{Vy}\PYG{o}{/}\PYG{n+nb+bp}{self}\PYG{o}{.}\PYG{n}{R} \PYG{o}{\PYGZhy{}} \PYG{n}{omega\PYGZus{}z}\PYG{o}{*}\PYG{p}{(}\PYG{n+nb+bp}{self}\PYG{o}{.}\PYG{n}{lx} \PYG{o}{+} \PYG{n+nb+bp}{self}\PYG{o}{.}\PYG{n}{ly}\PYG{p}{)}\PYG{o}{/}\PYG{p}{(}\PYG{l+m+mf}{2.0}\PYG{o}{*}\PYG{n+nb+bp}{self}\PYG{o}{.}\PYG{n}{R}\PYG{p}{)}\PYG{p}{,}
\PYG{n}{Vx}\PYG{o}{/}\PYG{n+nb+bp}{self}\PYG{o}{.}\PYG{n}{R} \PYG{o}{+} \PYG{n}{Vy}\PYG{o}{/}\PYG{n+nb+bp}{self}\PYG{o}{.}\PYG{n}{R} \PYG{o}{+} \PYG{n}{omega\PYGZus{}z}\PYG{o}{*}\PYG{p}{(}\PYG{n+nb+bp}{self}\PYG{o}{.}\PYG{n}{lx} \PYG{o}{+} \PYG{n+nb+bp}{self}\PYG{o}{.}\PYG{n}{ly}\PYG{p}{)}\PYG{o}{/}\PYG{p}{(}\PYG{l+m+mf}{2.0}\PYG{o}{*}\PYG{n+nb+bp}{self}\PYG{o}{.}\PYG{n}{R}\PYG{p}{)}\PYG{p}{,}
\PYG{n}{Vx}\PYG{o}{/}\PYG{n+nb+bp}{self}\PYG{o}{.}\PYG{n}{R} \PYG{o}{+} \PYG{n}{Vy}\PYG{o}{/}\PYG{n+nb+bp}{self}\PYG{o}{.}\PYG{n}{R} \PYG{o}{\PYGZhy{}} \PYG{n}{omega\PYGZus{}z}\PYG{o}{*}\PYG{p}{(}\PYG{n+nb+bp}{self}\PYG{o}{.}\PYG{n}{lx} \PYG{o}{+} \PYG{n+nb+bp}{self}\PYG{o}{.}\PYG{n}{ly}\PYG{p}{)}\PYG{o}{/}\PYG{p}{(}\PYG{l+m+mf}{2.0}\PYG{o}{*}\PYG{n+nb+bp}{self}\PYG{o}{.}\PYG{n}{R}\PYG{p}{)}\PYG{p}{,}
\PYG{n}{Vx}\PYG{o}{/}\PYG{n+nb+bp}{self}\PYG{o}{.}\PYG{n}{R} \PYG{o}{\PYGZhy{}} \PYG{n}{Vy}\PYG{o}{/}\PYG{n+nb+bp}{self}\PYG{o}{.}\PYG{n}{R} \PYG{o}{+} \PYG{n}{omega\PYGZus{}z}\PYG{o}{*}\PYG{p}{(}\PYG{n+nb+bp}{self}\PYG{o}{.}\PYG{n}{lx} \PYG{o}{+} \PYG{n+nb+bp}{self}\PYG{o}{.}\PYG{n}{ly}\PYG{p}{)}\PYG{o}{/}\PYG{p}{(}\PYG{l+m+mf}{2.0}\PYG{o}{*}\PYG{n+nb+bp}{self}\PYG{o}{.}\PYG{n}{R}\PYG{p}{)}
\PYG{p}{]}
\end{sphinxVerbatim}


\chapter{API Reference}
\label{\detokenize{api_reference:api-reference}}\label{\detokenize{api_reference::doc}}
\sphinxAtStartPar
Document the main interfaces, message formats, and public classes or services.
Add details for both app\sphinxhyphen{}side and robot\sphinxhyphen{}side APIs.


\section{Topics to cover}
\label{\detokenize{api_reference:topics-to-cover}}\begin{itemize}
\item {} 
\sphinxAtStartPar
Network messages and packet formats

\item {} 
\sphinxAtStartPar
ROS2 topics, services, and actions

\item {} 
\sphinxAtStartPar
Public classes and key methods

\item {} 
\sphinxAtStartPar
Configuration parameters and defaults

\end{itemize}


\chapter{Troubleshooting Guide}
\label{\detokenize{troubleshooting_guide:troubleshooting-guide}}\label{\detokenize{troubleshooting_guide::doc}}
\sphinxAtStartPar
List common issues, error messages, and recovery steps. Include diagnostic
commands and expected outputs where possible.


\section{Topics to cover}
\label{\detokenize{troubleshooting_guide:topics-to-cover}}\begin{itemize}
\item {} 
\sphinxAtStartPar
App connectivity and UDP issues

\item {} 
\sphinxAtStartPar
ROS2 node startup failures

\item {} 
\sphinxAtStartPar
Odometry or control anomalies

\item {} 
\sphinxAtStartPar
UI input or rendering problems

\item {} 
\sphinxAtStartPar
Log collection and debugging tips

\end{itemize}


\chapter{Appendix A: Object\sphinxhyphen{}Oriented Programming}
\label{\detokenize{appendix_a_object_oriented_programming:appendix-a-object-oriented-programming}}\label{\detokenize{appendix_a_object_oriented_programming::doc}}

\section{Overview}
\label{\detokenize{appendix_a_object_oriented_programming:overview}}
\sphinxAtStartPar
Object\sphinxhyphen{}oriented programming (OOP) models software as interacting objects that
encapsulate state and behavior. The goal is to make programs easier to reason
about by grouping related data and functions and enforcing clear boundaries.


\section{Core concepts}
\label{\detokenize{appendix_a_object_oriented_programming:core-concepts}}

\subsection{Classes and objects}
\label{\detokenize{appendix_a_object_oriented_programming:classes-and-objects}}\begin{itemize}
\item {} 
\sphinxAtStartPar
A class defines a blueprint for data (fields) and behavior (methods).

\item {} 
\sphinxAtStartPar
An object is an instance of a class created at runtime.

\end{itemize}


\subsection{Encapsulation}
\label{\detokenize{appendix_a_object_oriented_programming:encapsulation}}\begin{itemize}
\item {} 
\sphinxAtStartPar
Hide internal details and expose a minimal public interface.

\item {} 
\sphinxAtStartPar
Prevents accidental misuse and reduces coupling.

\end{itemize}


\subsection{Abstraction}
\label{\detokenize{appendix_a_object_oriented_programming:abstraction}}\begin{itemize}
\item {} 
\sphinxAtStartPar
Present only essential features and ignore implementation details.

\item {} 
\sphinxAtStartPar
Encourages focusing on what an object does, not how it does it.

\end{itemize}


\subsection{Inheritance}
\label{\detokenize{appendix_a_object_oriented_programming:inheritance}}\begin{itemize}
\item {} 
\sphinxAtStartPar
A subclass derives from a base class and can extend or override behavior.

\item {} 
\sphinxAtStartPar
Supports reuse but can introduce tight coupling if overused.

\end{itemize}


\subsection{Polymorphism}
\label{\detokenize{appendix_a_object_oriented_programming:polymorphism}}\begin{itemize}
\item {} 
\sphinxAtStartPar
Code can work with objects through a shared interface.

\item {} 
\sphinxAtStartPar
Allows different implementations to be swapped with minimal changes.

\end{itemize}


\section{Composition over inheritance}
\label{\detokenize{appendix_a_object_oriented_programming:composition-over-inheritance}}
\sphinxAtStartPar
Composition builds complex behavior by combining smaller objects, which helps
avoid fragile inheritance hierarchies. When possible, prefer composition and
interfaces to keep dependencies explicit.


\section{SOLID principles}
\label{\detokenize{appendix_a_object_oriented_programming:solid-principles}}\begin{itemize}
\item {} 
\sphinxAtStartPar
Single Responsibility: one reason to change.

\item {} 
\sphinxAtStartPar
Open/Closed: open for extension, closed for modification.

\item {} 
\sphinxAtStartPar
Liskov Substitution: subclasses must be usable where base classes are expected.

\item {} 
\sphinxAtStartPar
Interface Segregation: small, focused interfaces.

\item {} 
\sphinxAtStartPar
Dependency Inversion: depend on abstractions, not concrete types.

\end{itemize}


\section{Design patterns}
\label{\detokenize{appendix_a_object_oriented_programming:design-patterns}}
\sphinxAtStartPar
Common patterns that align with OOP:
\begin{itemize}
\item {} 
\sphinxAtStartPar
Strategy: swap algorithms without changing the caller.

\item {} 
\sphinxAtStartPar
Observer: notify subscribers on state changes.

\item {} 
\sphinxAtStartPar
Command: encapsulate an action as an object.

\item {} 
\sphinxAtStartPar
Factory: centralize object creation.

\end{itemize}


\section{Typical pitfalls}
\label{\detokenize{appendix_a_object_oriented_programming:typical-pitfalls}}\begin{itemize}
\item {} 
\sphinxAtStartPar
Deep inheritance hierarchies that obscure behavior.

\item {} 
\sphinxAtStartPar
God objects that violate single responsibility.

\item {} 
\sphinxAtStartPar
Exposing internal state directly, making invariants hard to enforce.

\end{itemize}


\section{Applied to robotics software}
\label{\detokenize{appendix_a_object_oriented_programming:applied-to-robotics-software}}
\sphinxAtStartPar
In robotics systems, OOP helps isolate concerns like communication, control,
and sensing. For example:
\begin{itemize}
\item {} 
\sphinxAtStartPar
A network manager can encapsulate UDP send/receive behavior.

\item {} 
\sphinxAtStartPar
A controller class can wrap control logic and state updates.

\item {} 
\sphinxAtStartPar
A UI layer can treat the model as a stable interface and remain decoupled
from data sources.

\end{itemize}


\section{Example: interface\sphinxhyphen{}driven control}
\label{\detokenize{appendix_a_object_oriented_programming:example-interface-driven-control}}
\begin{sphinxVerbatim}[commandchars=\\\{\}]
\PYG{k}{class}\PYG{+w}{ }\PYG{n+nc}{VelocityController}\PYG{p}{:}
    \PYG{k}{def}\PYG{+w}{ }\PYG{n+nf}{compute}\PYG{p}{(}\PYG{n+nb+bp}{self}\PYG{p}{,} \PYG{n}{target}\PYG{p}{,} \PYG{n}{state}\PYG{p}{)}\PYG{p}{:}
        \PYG{k}{raise} \PYG{n+ne}{NotImplementedError}

\PYG{k}{class}\PYG{+w}{ }\PYG{n+nc}{PDController}\PYG{p}{(}\PYG{n}{VelocityController}\PYG{p}{)}\PYG{p}{:}
    \PYG{k}{def}\PYG{+w}{ }\PYG{n+nf}{compute}\PYG{p}{(}\PYG{n+nb+bp}{self}\PYG{p}{,} \PYG{n}{target}\PYG{p}{,} \PYG{n}{state}\PYG{p}{)}\PYG{p}{:}
        \PYG{c+c1}{\PYGZsh{} PD math here}
        \PYG{k}{return} \PYG{n}{vx}\PYG{p}{,} \PYG{n}{vy}\PYG{p}{,} \PYG{n}{wz}

\PYG{k}{def}\PYG{+w}{ }\PYG{n+nf}{step}\PYG{p}{(}\PYG{n}{controller}\PYG{p}{,} \PYG{n}{target}\PYG{p}{,} \PYG{n}{state}\PYG{p}{)}\PYG{p}{:}
    \PYG{k}{return} \PYG{n}{controller}\PYG{o}{.}\PYG{n}{compute}\PYG{p}{(}\PYG{n}{target}\PYG{p}{,} \PYG{n}{state}\PYG{p}{)}
\end{sphinxVerbatim}

\sphinxAtStartPar
This keeps the control pipeline stable while allowing multiple controllers to
be used during testing or future upgrades.


\chapter{Appendix B: MVC Architecture}
\label{\detokenize{appendix_b_mvc_architecture:appendix-b-mvc-architecture}}\label{\detokenize{appendix_b_mvc_architecture::doc}}

\section{Overview}
\label{\detokenize{appendix_b_mvc_architecture:overview}}
\sphinxAtStartPar
Model\sphinxhyphen{}View\sphinxhyphen{}Controller (MVC) separates data, presentation, and user input into
three roles. This separation improves maintainability by reducing coupling.


\section{Roles}
\label{\detokenize{appendix_b_mvc_architecture:roles}}

\subsection{Model}
\label{\detokenize{appendix_b_mvc_architecture:model}}\begin{itemize}
\item {} 
\sphinxAtStartPar
Owns the application state and business rules.

\item {} 
\sphinxAtStartPar
Exposes operations that modify data in a controlled way.

\end{itemize}


\subsection{View}
\label{\detokenize{appendix_b_mvc_architecture:view}}\begin{itemize}
\item {} 
\sphinxAtStartPar
Renders the model state to the user.

\item {} 
\sphinxAtStartPar
Should be mostly passive and not contain business logic.

\end{itemize}


\subsection{Controller}
\label{\detokenize{appendix_b_mvc_architecture:controller}}\begin{itemize}
\item {} 
\sphinxAtStartPar
Interprets user input and updates the model.

\item {} 
\sphinxAtStartPar
Chooses which view to refresh or navigate to.

\end{itemize}


\section{Typical flow}
\label{\detokenize{appendix_b_mvc_architecture:typical-flow}}\begin{enumerate}
\sphinxsetlistlabels{\arabic}{enumi}{enumii}{}{.}%
\item {} 
\sphinxAtStartPar
User interacts with the View.

\item {} 
\sphinxAtStartPar
Controller translates the action into model updates.

\item {} 
\sphinxAtStartPar
Model updates its state and notifies views.

\item {} 
\sphinxAtStartPar
View re\sphinxhyphen{}renders using the new model state.

\end{enumerate}


\section{Benefits}
\label{\detokenize{appendix_b_mvc_architecture:benefits}}\begin{itemize}
\item {} 
\sphinxAtStartPar
Clear separation of responsibilities.

\item {} 
\sphinxAtStartPar
Easier testing of logic by isolating the model.

\item {} 
\sphinxAtStartPar
Views can be swapped without changing the model.

\end{itemize}


\section{Tradeoffs}
\label{\detokenize{appendix_b_mvc_architecture:tradeoffs}}\begin{itemize}
\item {} 
\sphinxAtStartPar
Controllers can grow large if not structured.

\item {} 
\sphinxAtStartPar
Requires discipline to keep business logic in the model.

\end{itemize}


\section{Applied to a robotics UI}
\label{\detokenize{appendix_b_mvc_architecture:applied-to-a-robotics-ui}}
\sphinxAtStartPar
In a waypoint\sphinxhyphen{}setting UI:
\begin{itemize}
\item {} 
\sphinxAtStartPar
Model holds the waypoint list, active target, and plan ID.

\item {} 
\sphinxAtStartPar
View renders the field, robot position, and waypoint markers.

\item {} 
\sphinxAtStartPar
Controller translates touch events into model updates and network sends.

\end{itemize}


\section{Implementation hint}
\label{\detokenize{appendix_b_mvc_architecture:implementation-hint}}
\sphinxAtStartPar
Use small controllers to handle independent features (e.g., waypoint editing,
return\sphinxhyphen{}to\sphinxhyphen{}base). This avoids a single controller that becomes a bottleneck.


\chapter{Appendix C: MVVM Architecture}
\label{\detokenize{appendix_c_mvvm_architecture:appendix-c-mvvm-architecture}}\label{\detokenize{appendix_c_mvvm_architecture::doc}}

\section{Overview}
\label{\detokenize{appendix_c_mvvm_architecture:overview}}
\sphinxAtStartPar
Model\sphinxhyphen{}View\sphinxhyphen{}ViewModel (MVVM) emphasizes data binding between view and view\sphinxhyphen{}model.
It is common in UI frameworks with declarative views.


\section{Roles}
\label{\detokenize{appendix_c_mvvm_architecture:roles}}

\subsection{Model}
\label{\detokenize{appendix_c_mvvm_architecture:model}}\begin{itemize}
\item {} 
\sphinxAtStartPar
Domain data and business rules.

\item {} 
\sphinxAtStartPar
Independent of UI technologies.

\end{itemize}


\subsection{View}
\label{\detokenize{appendix_c_mvvm_architecture:view}}\begin{itemize}
\item {} 
\sphinxAtStartPar
UI elements and layout.

\item {} 
\sphinxAtStartPar
Binds directly to view\sphinxhyphen{}model properties.

\end{itemize}


\subsection{ViewModel}
\label{\detokenize{appendix_c_mvvm_architecture:viewmodel}}\begin{itemize}
\item {} 
\sphinxAtStartPar
Adapts model data into UI\sphinxhyphen{}friendly formats.

\item {} 
\sphinxAtStartPar
Exposes commands or actions that the view can invoke.

\end{itemize}


\section{Data binding}
\label{\detokenize{appendix_c_mvvm_architecture:data-binding}}
\sphinxAtStartPar
Bindings update the UI when model data changes, and optionally update the model
when the user edits the UI. This reduces glue code but requires careful state
management.


\section{Benefits}
\label{\detokenize{appendix_c_mvvm_architecture:benefits}}\begin{itemize}
\item {} 
\sphinxAtStartPar
Strong separation between view and logic.

\item {} 
\sphinxAtStartPar
Easier to test view\sphinxhyphen{}model behavior.

\item {} 
\sphinxAtStartPar
Reduced UI boilerplate.

\end{itemize}


\section{Tradeoffs}
\label{\detokenize{appendix_c_mvvm_architecture:tradeoffs}}\begin{itemize}
\item {} 
\sphinxAtStartPar
Overuse of bindings can make data flow hard to trace.

\item {} 
\sphinxAtStartPar
View\sphinxhyphen{}models can become too large without clear boundaries.

\end{itemize}


\section{Applied to target setting}
\label{\detokenize{appendix_c_mvvm_architecture:applied-to-target-setting}}
\sphinxAtStartPar
A view\sphinxhyphen{}model might expose:
\begin{itemize}
\item {} 
\sphinxAtStartPar
A list of waypoints as UI\sphinxhyphen{}friendly points.

\item {} 
\sphinxAtStartPar
A selected waypoint index.

\item {} 
\sphinxAtStartPar
Commands for send, edit, and return actions.

\end{itemize}

\sphinxAtStartPar
This allows the view to remain simple and focused on rendering and input.


\chapter{Appendix D: ROS2 Overview}
\label{\detokenize{appendix_d_ros2_overview:appendix-d-ros2-overview}}\label{\detokenize{appendix_d_ros2_overview::doc}}

\section{What is ROS2}
\label{\detokenize{appendix_d_ros2_overview:what-is-ros2}}
\sphinxAtStartPar
ROS2 (Robot Operating System 2) is a middleware framework for building robot
applications. It provides communication, tools, and libraries for composing
complex robot systems from modular nodes.


\section{Key concepts}
\label{\detokenize{appendix_d_ros2_overview:key-concepts}}

\subsection{Nodes}
\label{\detokenize{appendix_d_ros2_overview:nodes}}\begin{itemize}
\item {} 
\sphinxAtStartPar
Independent processes that perform specific tasks.

\item {} 
\sphinxAtStartPar
Nodes communicate via topics, services, and actions.

\end{itemize}


\subsection{Topics}
\label{\detokenize{appendix_d_ros2_overview:topics}}\begin{itemize}
\item {} 
\sphinxAtStartPar
Publish/subscribe channels for streaming data.

\item {} 
\sphinxAtStartPar
Designed for high\sphinxhyphen{}frequency sensor and state updates.

\end{itemize}


\subsection{Services}
\label{\detokenize{appendix_d_ros2_overview:services}}\begin{itemize}
\item {} 
\sphinxAtStartPar
Synchronous request/response interactions.

\item {} 
\sphinxAtStartPar
Used for operations like configuration or queries.

\end{itemize}


\subsection{Actions}
\label{\detokenize{appendix_d_ros2_overview:actions}}\begin{itemize}
\item {} 
\sphinxAtStartPar
Long\sphinxhyphen{}running goals with feedback and cancellation.

\item {} 
\sphinxAtStartPar
Suitable for navigation and planning tasks.

\end{itemize}


\subsection{Parameters}
\label{\detokenize{appendix_d_ros2_overview:parameters}}\begin{itemize}
\item {} 
\sphinxAtStartPar
Runtime configuration values stored on nodes.

\item {} 
\sphinxAtStartPar
Can be read and set externally.

\end{itemize}


\section{QoS (Quality of Service)}
\label{\detokenize{appendix_d_ros2_overview:qos-quality-of-service}}
\sphinxAtStartPar
ROS2 exposes DDS QoS settings like reliability, durability, and history depth.
Selecting the right QoS is critical in lossy networks or real\sphinxhyphen{}time systems.


\section{Packages and workspaces}
\label{\detokenize{appendix_d_ros2_overview:packages-and-workspaces}}\begin{itemize}
\item {} 
\sphinxAtStartPar
A package is a unit of build and deployment.

\item {} 
\sphinxAtStartPar
A workspace contains multiple packages and their build artifacts.

\end{itemize}


\section{Lifecycle management}
\label{\detokenize{appendix_d_ros2_overview:lifecycle-management}}
\sphinxAtStartPar
ROS2 supports managed nodes with lifecycle states (unconfigured, inactive,
active, finalized). This helps bring up systems in a predictable order.


\section{Why ROS2 for robotics}
\label{\detokenize{appendix_d_ros2_overview:why-ros2-for-robotics}}\begin{itemize}
\item {} 
\sphinxAtStartPar
Scalable communication infrastructure.

\item {} 
\sphinxAtStartPar
Tooling for visualization and debugging.

\item {} 
\sphinxAtStartPar
Large ecosystem of drivers and algorithms.

\end{itemize}


\chapter{Appendix E: Distributed Systems}
\label{\detokenize{appendix_e_distributed_systems:appendix-e-distributed-systems}}\label{\detokenize{appendix_e_distributed_systems::doc}}

\section{Overview}
\label{\detokenize{appendix_e_distributed_systems:overview}}
\sphinxAtStartPar
A distributed system is a collection of independent processes that cooperate
via a network. Robotics stacks are inherently distributed, with sensors,
controllers, and user interfaces running on separate nodes.


\section{Key challenges}
\label{\detokenize{appendix_e_distributed_systems:key-challenges}}

\subsection{Latency and jitter}
\label{\detokenize{appendix_e_distributed_systems:latency-and-jitter}}\begin{itemize}
\item {} 
\sphinxAtStartPar
Network delays vary over time.

\item {} 
\sphinxAtStartPar
Control loops must tolerate nondeterministic timing.

\end{itemize}


\subsection{Partial failures}
\label{\detokenize{appendix_e_distributed_systems:partial-failures}}\begin{itemize}
\item {} 
\sphinxAtStartPar
Components can fail independently.

\item {} 
\sphinxAtStartPar
Systems should degrade gracefully when nodes disconnect.

\end{itemize}


\subsection{Consistency}
\label{\detokenize{appendix_e_distributed_systems:consistency}}\begin{itemize}
\item {} 
\sphinxAtStartPar
Shared state is hard to keep synchronized.

\item {} 
\sphinxAtStartPar
Stale data can lead to incorrect decisions.

\end{itemize}


\section{Time and ordering}
\label{\detokenize{appendix_e_distributed_systems:time-and-ordering}}\begin{itemize}
\item {} 
\sphinxAtStartPar
Clocks drift between machines.

\item {} 
\sphinxAtStartPar
Use timestamps and sequence numbers to detect ordering issues.

\end{itemize}


\section{CAP tradeoff}
\label{\detokenize{appendix_e_distributed_systems:cap-tradeoff}}
\sphinxAtStartPar
The CAP theorem highlights tradeoffs between consistency, availability, and
partition tolerance. Robotics systems often prioritize availability and timely
responses over strict consistency.


\section{Patterns}
\label{\detokenize{appendix_e_distributed_systems:patterns}}\begin{itemize}
\item {} 
\sphinxAtStartPar
Publish/subscribe for streaming state.

\item {} 
\sphinxAtStartPar
Request/response for configuration.

\item {} 
\sphinxAtStartPar
Heartbeats for liveness detection.

\end{itemize}


\section{Applied to this project}
\label{\detokenize{appendix_e_distributed_systems:applied-to-this-project}}
\sphinxAtStartPar
Waypoint commands and odometry updates are exchanged over UDP, which trades
reliability for low latency. The app must tolerate packet loss and occasional
stale data while keeping the UI responsive.


\chapter{Appendix F: Monolithic Architecture}
\label{\detokenize{appendix_f_monolithic_architecture:appendix-f-monolithic-architecture}}\label{\detokenize{appendix_f_monolithic_architecture::doc}}

\section{Overview}
\label{\detokenize{appendix_f_monolithic_architecture:overview}}
\sphinxAtStartPar
A monolithic architecture bundles application components into a single deployable
unit. This contrasts with microservices or distributed components.


\section{Benefits}
\label{\detokenize{appendix_f_monolithic_architecture:benefits}}\begin{itemize}
\item {} 
\sphinxAtStartPar
Simple deployment and debugging.

\item {} 
\sphinxAtStartPar
No network overhead between internal modules.

\item {} 
\sphinxAtStartPar
Easier to maintain transaction boundaries.

\end{itemize}


\section{Drawbacks}
\label{\detokenize{appendix_f_monolithic_architecture:drawbacks}}\begin{itemize}
\item {} 
\sphinxAtStartPar
Tightly coupled components reduce flexibility.

\item {} 
\sphinxAtStartPar
Scaling requires duplicating the entire system.

\item {} 
\sphinxAtStartPar
Large codebases can slow development without strong modularity.

\end{itemize}


\section{Modular monolith}
\label{\detokenize{appendix_f_monolithic_architecture:modular-monolith}}
\sphinxAtStartPar
A modular monolith preserves a single deployment but enforces internal module
boundaries. This approach often provides the best of both worlds when the
system does not require independent scaling.


\section{Applied to robotics tooling}
\label{\detokenize{appendix_f_monolithic_architecture:applied-to-robotics-tooling}}
\sphinxAtStartPar
A standalone UI app can be treated as a monolith with clear internal modules
for networking, input handling, rendering, and data persistence. This keeps
deployment simple while still enabling maintainability.


\chapter{Appendix G: ROS2 Architecture}
\label{\detokenize{appendix_g_ros2_architecture:appendix-g-ros2-architecture}}\label{\detokenize{appendix_g_ros2_architecture::doc}}

\section{Layered structure}
\label{\detokenize{appendix_g_ros2_architecture:layered-structure}}
\sphinxAtStartPar
ROS2 is built as a stack of layers:
\begin{itemize}
\item {} 
\sphinxAtStartPar
Client libraries: rclcpp (C++), rclpy (Python)

\item {} 
\sphinxAtStartPar
ROS client library: rcl

\item {} 
\sphinxAtStartPar
Middleware interface: rmw

\item {} 
\sphinxAtStartPar
DDS implementation: Fast DDS, Cyclone DDS, or Connext

\end{itemize}


\section{Executors}
\label{\detokenize{appendix_g_ros2_architecture:executors}}
\sphinxAtStartPar
Executors schedule callbacks for subscriptions, timers, and services. Common
executors include single\sphinxhyphen{}threaded and multi\sphinxhyphen{}threaded variants. Callback groups
control concurrency and prevent data races.


\section{Composition}
\label{\detokenize{appendix_g_ros2_architecture:composition}}
\sphinxAtStartPar
ROS2 supports composable nodes that run in a single process. This reduces
overhead and improves performance when nodes are tightly coupled.


\section{Lifecycle nodes}
\label{\detokenize{appendix_g_ros2_architecture:lifecycle-nodes}}
\sphinxAtStartPar
Lifecycle nodes provide deterministic startup and shutdown behavior. They move
through states like unconfigured, inactive, active, and finalized.


\section{DDS details}
\label{\detokenize{appendix_g_ros2_architecture:dds-details}}
\sphinxAtStartPar
DDS handles discovery, serialization, and transport. QoS policies determine how
data is delivered and what guarantees are provided.


\section{Best practices}
\label{\detokenize{appendix_g_ros2_architecture:best-practices}}\begin{itemize}
\item {} 
\sphinxAtStartPar
Choose QoS profiles based on network conditions.

\item {} 
\sphinxAtStartPar
Use namespaces to keep large systems organized.

\item {} 
\sphinxAtStartPar
Keep topics small and publish at stable rates.

\end{itemize}


\chapter{Appendix H: UDP}
\label{\detokenize{appendix_h_udp:appendix-h-udp}}\label{\detokenize{appendix_h_udp::doc}}

\section{Overview}
\label{\detokenize{appendix_h_udp:overview}}
\sphinxAtStartPar
User Datagram Protocol (UDP) is a connectionless transport protocol that sends
independent datagrams. It is fast and lightweight but provides no delivery or
ordering guarantees.


\section{Characteristics}
\label{\detokenize{appendix_h_udp:characteristics}}\begin{itemize}
\item {} 
\sphinxAtStartPar
Connectionless and stateless.

\item {} 
\sphinxAtStartPar
Packets can be dropped, duplicated, or reordered.

\item {} 
\sphinxAtStartPar
Lower overhead than TCP.

\end{itemize}


\section{Packet size and MTU}
\label{\detokenize{appendix_h_udp:packet-size-and-mtu}}
\sphinxAtStartPar
UDP payloads that exceed the path MTU may be fragmented, which increases loss
risk. Design packet formats to be compact and avoid fragmentation where
possible.


\section{Checksums}
\label{\detokenize{appendix_h_udp:checksums}}
\sphinxAtStartPar
UDP includes a checksum for error detection but no retransmission. Applications
must decide whether to drop or tolerate corrupted packets.


\section{Best practices for control systems}
\label{\detokenize{appendix_h_udp:best-practices-for-control-systems}}\begin{itemize}
\item {} 
\sphinxAtStartPar
Include sequence numbers or plan IDs to identify stale data.

\item {} 
\sphinxAtStartPar
Use fixed\sphinxhyphen{}size binary formats for consistent parsing.

\item {} 
\sphinxAtStartPar
Keep send rates bounded to avoid congestion.

\item {} 
\sphinxAtStartPar
For critical commands, consider sending multiple times or adding acknowledgments.

\end{itemize}


\section{Applied to target setting}
\label{\detokenize{appendix_h_udp:applied-to-target-setting}}
\sphinxAtStartPar
The waypoint protocol favors low latency and simplicity. The app can tolerate
occasional packet loss while maintaining responsiveness.


\chapter{Appendix I: Quaternions}
\label{\detokenize{appendix_i_quaternions:appendix-i-quaternions}}\label{\detokenize{appendix_i_quaternions::doc}}

\section{Overview}
\label{\detokenize{appendix_i_quaternions:overview}}
\sphinxAtStartPar
Quaternions represent 3D rotations without gimbal lock. A quaternion is written
as \$q = (w, x, y, z)\$ where \$w\$ is the scalar component and \$(x, y, z)\$ is the
vector component.


\section{Unit quaternions}
\label{\detokenize{appendix_i_quaternions:unit-quaternions}}
\sphinxAtStartPar
A rotation quaternion is unit length:
\begin{equation*}
\begin{split}\|q\| = \sqrt{w^2 + x^2 + y^2 + z^2} = 1\end{split}
\end{equation*}

\section{Quaternion multiplication}
\label{\detokenize{appendix_i_quaternions:quaternion-multiplication}}
\sphinxAtStartPar
Combining rotations uses quaternion multiplication, which is not commutative.
If \$q\_1\$ and \$q\_2\$ are rotations, the composed rotation is \$q = q\_2 otimes q\_1\$.


\section{Axis\sphinxhyphen{}angle conversion}
\label{\detokenize{appendix_i_quaternions:axis-angle-conversion}}
\sphinxAtStartPar
For rotation angle \$theta\$ around unit axis \$hat\{u\}\$:
\begin{equation*}
\begin{split}q = \left(\cos(\theta/2),\ \hat{u}_x\sin(\theta/2),\ \hat{u}_y\sin(\theta/2),\ \hat{u}_z\sin(\theta/2)\right)\end{split}
\end{equation*}

\section{Yaw extraction (2D robots)}
\label{\detokenize{appendix_i_quaternions:yaw-extraction-2d-robots}}
\sphinxAtStartPar
For planar motion, yaw can be extracted from a quaternion:
\begin{equation*}
\begin{split}\text{yaw} = \arctan2(2(wz + xy),\ 1 - 2(y^2 + z^2))\end{split}
\end{equation*}

\section{Numerical stability}
\label{\detokenize{appendix_i_quaternions:numerical-stability}}\begin{itemize}
\item {} 
\sphinxAtStartPar
Normalize quaternions after integration.

\item {} 
\sphinxAtStartPar
Avoid repeated small\sphinxhyphen{}angle conversions without renormalization.

\end{itemize}


\section{Applied to odometry}
\label{\detokenize{appendix_i_quaternions:applied-to-odometry}}
\sphinxAtStartPar
Yaw is computed from quaternion orientation and then used in frame transforms
for velocities and positions.


\chapter{Appendix J: Kinematics}
\label{\detokenize{appendix_j_kinematics:appendix-j-kinematics}}\label{\detokenize{appendix_j_kinematics::doc}}

\section{Overview}
\label{\detokenize{appendix_j_kinematics:overview}}
\sphinxAtStartPar
Kinematics describes the relationship between motion variables without
considering forces. For robots, this links wheel velocities to body motion.


\section{Forward kinematics}
\label{\detokenize{appendix_j_kinematics:forward-kinematics}}
\sphinxAtStartPar
Given wheel speeds, compute body velocities. For a holonomic robot, the
relationship can be expressed as:
\begin{equation*}
\begin{split}\mathbf{v} = \mathbf{J} \cdot \mathbf{\omega}\end{split}
\end{equation*}
\sphinxAtStartPar
where \$mathbf\{J\}\$ is the kinematic matrix and \$mathbf\{omega\}\$ are wheel rates.


\section{Inverse kinematics}
\label{\detokenize{appendix_j_kinematics:inverse-kinematics}}
\sphinxAtStartPar
Given desired body velocities, compute wheel speeds:
\begin{equation*}
\begin{split}\mathbf{\omega} = \mathbf{J}^{-1} \cdot \mathbf{v}\end{split}
\end{equation*}
\sphinxAtStartPar
If \$mathbf\{J\}\$ is not square, a pseudoinverse is used.


\section{Coordinate frames}
\label{\detokenize{appendix_j_kinematics:coordinate-frames}}\begin{itemize}
\item {} 
\sphinxAtStartPar
Local frame: attached to the robot body.

\item {} 
\sphinxAtStartPar
Global frame: fixed to the environment or map.

\end{itemize}

\sphinxAtStartPar
Transform between frames using rotation matrices derived from yaw.


\section{Discretization}
\label{\detokenize{appendix_j_kinematics:discretization}}
\sphinxAtStartPar
Velocities are integrated over time to compute position:
\begin{equation*}
\begin{split}x_{k+1} = x_k + v_x \Delta t\end{split}
\end{equation*}
\sphinxAtStartPar
and similarly for \$y\$ and \$theta\$.


\section{Applied to this system}
\label{\detokenize{appendix_j_kinematics:applied-to-this-system}}
\sphinxAtStartPar
The kinematic model converts desired planar velocities into wheel velocities
and back into odometry updates.


\chapter{Appendix K: State Machine}
\label{\detokenize{appendix_k_state_machine:appendix-k-state-machine}}\label{\detokenize{appendix_k_state_machine::doc}}

\section{Overview}
\label{\detokenize{appendix_k_state_machine:overview}}
\sphinxAtStartPar
A state machine models a system with a finite number of modes. Each state has
well\sphinxhyphen{}defined behavior and transitions that occur on events or conditions.


\section{Key terms}
\label{\detokenize{appendix_k_state_machine:key-terms}}\begin{itemize}
\item {} 
\sphinxAtStartPar
State: a named mode of operation.

\item {} 
\sphinxAtStartPar
Transition: a change from one state to another.

\item {} 
\sphinxAtStartPar
Guard: a condition that must be true for a transition.

\item {} 
\sphinxAtStartPar
Action: logic executed on entry, exit, or during a state.

\end{itemize}


\section{Design guidelines}
\label{\detokenize{appendix_k_state_machine:design-guidelines}}\begin{itemize}
\item {} 
\sphinxAtStartPar
Keep states orthogonal: one reason to change.

\item {} 
\sphinxAtStartPar
Prefer explicit transitions over implicit side effects.

\item {} 
\sphinxAtStartPar
Log transitions for debugging.

\end{itemize}


\section{Hierarchical state machines}
\label{\detokenize{appendix_k_state_machine:hierarchical-state-machines}}
\sphinxAtStartPar
Nested states allow shared behavior at higher levels. This reduces duplication
when multiple states share common logic.


\section{Applied to navigation}
\label{\detokenize{appendix_k_state_machine:applied-to-navigation}}
\sphinxAtStartPar
The navigation pipeline can be modeled as:
\begin{itemize}
\item {} 
\sphinxAtStartPar
IDLE: waiting for a waypoint.

\item {} 
\sphinxAtStartPar
NAVIGATING: moving toward a target.

\item {} 
\sphinxAtStartPar
EDIT: applying a waypoint update.

\item {} 
\sphinxAtStartPar
PAUSED: holding position.

\item {} 
\sphinxAtStartPar
RETURNING: going back to a visited waypoint.

\end{itemize}

\sphinxAtStartPar
This structure simplifies reasoning about behavior and failure handling.


\chapter{Appendix L: PID Control}
\label{\detokenize{appendix_l_pid_control:appendix-l-pid-control}}\label{\detokenize{appendix_l_pid_control::doc}}

\section{Overview}
\label{\detokenize{appendix_l_pid_control:overview}}
\sphinxAtStartPar
PID control combines proportional, integral, and derivative terms to reduce
error between a desired value and the measured value.


\section{PID equation}
\label{\detokenize{appendix_l_pid_control:pid-equation}}\begin{equation*}
\begin{split}u(t) = K_p e(t) + K_i \int_0^t e(\tau) d\tau + K_d \frac{de(t)}{dt}\end{split}
\end{equation*}

\section{Discrete form}
\label{\detokenize{appendix_l_pid_control:discrete-form}}\begin{equation*}
\begin{split}u_k = K_p e_k + K_i \sum_{i=0}^k e_i \Delta t + K_d \frac{e_k - e_{k-1}}{\Delta t}\end{split}
\end{equation*}

\section{Tuning}
\label{\detokenize{appendix_l_pid_control:tuning}}\begin{itemize}
\item {} 
\sphinxAtStartPar
\$K\_p\$ sets responsiveness but can overshoot.

\item {} 
\sphinxAtStartPar
\$K\_i\$ removes steady\sphinxhyphen{}state error but can wind up.

\item {} 
\sphinxAtStartPar
\$K\_d\$ dampens oscillation but amplifies noise.

\end{itemize}


\section{Anti\sphinxhyphen{}windup}
\label{\detokenize{appendix_l_pid_control:anti-windup}}
\sphinxAtStartPar
Clamp the integral term when actuators saturate:
\begin{equation*}
\begin{split}I_k = \text{clamp}(I_k, I_{min}, I_{max})\end{split}
\end{equation*}

\section{Filtering}
\label{\detokenize{appendix_l_pid_control:filtering}}
\sphinxAtStartPar
Derivative terms often use low\sphinxhyphen{}pass filtering to reduce noise.


\section{Applied to robot motion}
\label{\detokenize{appendix_l_pid_control:applied-to-robot-motion}}
\sphinxAtStartPar
The controller uses PD components on position and yaw with filtered derivative
terms for smoother navigation.


\chapter{Appendix M: Localization}
\label{\detokenize{appendix_m_localization:appendix-m-localization}}\label{\detokenize{appendix_m_localization::doc}}

\section{Overview}
\label{\detokenize{appendix_m_localization:overview}}
\sphinxAtStartPar
Localization estimates a robot’s position and orientation in a global frame.
It combines motion models with sensor measurements.


\section{Sources of localization}
\label{\detokenize{appendix_m_localization:sources-of-localization}}\begin{itemize}
\item {} 
\sphinxAtStartPar
Odometry: integrates wheel or motor motion.

\item {} 
\sphinxAtStartPar
IMU: measures angular velocity and acceleration.

\item {} 
\sphinxAtStartPar
Visual markers or lidar: provide absolute corrections.

\end{itemize}


\section{Probabilistic view}
\label{\detokenize{appendix_m_localization:probabilistic-view}}
\sphinxAtStartPar
State estimation often uses a probabilistic model:
\begin{equation*}
\begin{split}x_k = f(x_{k-1}, u_k) + w_k
z_k = h(x_k) + v_k\end{split}
\end{equation*}
\sphinxAtStartPar
where \$w\_k\$ and \$v\_k\$ represent process and measurement noise.


\section{Filters}
\label{\detokenize{appendix_m_localization:filters}}\begin{itemize}
\item {} 
\sphinxAtStartPar
Extended Kalman Filter (EKF): linearizes nonlinear models.

\item {} 
\sphinxAtStartPar
Particle Filter: represents state as particles.

\end{itemize}


\section{Frames}
\label{\detokenize{appendix_m_localization:frames}}
\sphinxAtStartPar
Common frames in robotics include:
\begin{itemize}
\item {} 
\sphinxAtStartPar
base\_link: robot body.

\item {} 
\sphinxAtStartPar
odom: local frame for short\sphinxhyphen{}term accuracy.

\item {} 
\sphinxAtStartPar
map: global reference frame.

\end{itemize}


\section{Applied to this system}
\label{\detokenize{appendix_m_localization:applied-to-this-system}}
\sphinxAtStartPar
Odometry provides short\sphinxhyphen{}term pose estimates. For long\sphinxhyphen{}term accuracy, periodic
corrections can be applied using external references.


\chapter{Appendix N: Odometry}
\label{\detokenize{appendix_n_odometry:appendix-n-odometry}}\label{\detokenize{appendix_n_odometry::doc}}

\section{Overview}
\label{\detokenize{appendix_n_odometry:overview}}
\sphinxAtStartPar
Odometry estimates robot motion by integrating velocity over time. It is fast
and local but accumulates drift due to wheel slip and noise.


\section{Kinematic integration}
\label{\detokenize{appendix_n_odometry:kinematic-integration}}
\sphinxAtStartPar
Given linear velocities in the global frame:
\begin{equation*}
\begin{split}x_{k+1} = x_k + v_x \Delta t
y_{k+1} = y_k + v_y \Delta t\end{split}
\end{equation*}
\sphinxAtStartPar
Yaw integration uses angular velocity:
\begin{equation*}
\begin{split}\theta_{k+1} = \theta_k + \omega_z \Delta t\end{split}
\end{equation*}

\section{Error sources}
\label{\detokenize{appendix_n_odometry:error-sources}}\begin{itemize}
\item {} 
\sphinxAtStartPar
Wheel slip or uneven terrain.

\item {} 
\sphinxAtStartPar
Encoder quantization.

\item {} 
\sphinxAtStartPar
Incorrect wheel radius or geometry parameters.

\end{itemize}


\section{Mitigation}
\label{\detokenize{appendix_n_odometry:mitigation}}\begin{itemize}
\item {} 
\sphinxAtStartPar
Periodic calibration of wheel parameters.

\item {} 
\sphinxAtStartPar
Fusing odometry with IMU or external sensors.

\item {} 
\sphinxAtStartPar
Avoiding long open\sphinxhyphen{}loop runs without correction.

\end{itemize}


\section{Applied to this project}
\label{\detokenize{appendix_n_odometry:applied-to-this-project}}
\sphinxAtStartPar
Odometry is derived from motor velocities and transformed into the global frame
for display and control.


\chapter{Appendix O: Waypoint Navigation}
\label{\detokenize{appendix_o_waypoint_navigation:appendix-o-waypoint-navigation}}\label{\detokenize{appendix_o_waypoint_navigation::doc}}

\section{Overview}
\label{\detokenize{appendix_o_waypoint_navigation:overview}}
\sphinxAtStartPar
Waypoint navigation drives a robot through a sequence of target positions.
Key components include path planning, control, and arrival detection.


\section{Path planning}
\label{\detokenize{appendix_o_waypoint_navigation:path-planning}}\begin{itemize}
\item {} 
\sphinxAtStartPar
For sparse waypoints, straight\sphinxhyphen{}line segments are common.

\item {} 
\sphinxAtStartPar
For smoother trajectories, use splines or polynomial paths.

\end{itemize}


\section{Arrival detection}
\label{\detokenize{appendix_o_waypoint_navigation:arrival-detection}}
\sphinxAtStartPar
A waypoint is considered reached when distance and timing constraints are met:
\begin{equation*}
\begin{split}\text{dist} = \sqrt{(x_t - x)^2 + (y_t - y)^2}\end{split}
\end{equation*}
\sphinxAtStartPar
If dist is below a tolerance for a specified duration, the waypoint is complete.


\section{Handling edits}
\label{\detokenize{appendix_o_waypoint_navigation:handling-edits}}
\sphinxAtStartPar
When waypoints are edited, smoothing can reduce abrupt changes in command
velocity. A simple low\sphinxhyphen{}pass update is:
\begin{equation*}
\begin{split}x_{new} = (1 - \alpha) x_{old} + \alpha x_{update}\end{split}
\end{equation*}

\section{Return behavior}
\label{\detokenize{appendix_o_waypoint_navigation:return-behavior}}
\sphinxAtStartPar
Visited waypoints can be stored in a stack to support a return\sphinxhyphen{}to\sphinxhyphen{}base function.


\section{Applied to this system}
\label{\detokenize{appendix_o_waypoint_navigation:applied-to-this-system}}
\sphinxAtStartPar
The app sends waypoint lists and updates to the robot. The controller processes
the queue, updates the active target, and transitions between states.


\chapter{Appendix P: Complementary Filtering}
\label{\detokenize{appendix_p_complementary_filtering:appendix-p-complementary-filtering}}\label{\detokenize{appendix_p_complementary_filtering::doc}}

\section{Overview}
\label{\detokenize{appendix_p_complementary_filtering:overview}}
\sphinxAtStartPar
A complementary filter combines two signals with different frequency qualities.
One signal provides good low\sphinxhyphen{}frequency behavior, the other provides high\sphinxhyphen{}frequency
response. The sum provides a balanced estimate.


\section{Basic equation}
\label{\detokenize{appendix_p_complementary_filtering:basic-equation}}\begin{equation*}
\begin{split}y = \alpha y_{fast} + (1 - \alpha) y_{slow}\end{split}
\end{equation*}
\sphinxAtStartPar
where \$0 \textless{} alpha \textless{} 1\$.


\section{Why it works}
\label{\detokenize{appendix_p_complementary_filtering:why-it-works}}\begin{itemize}
\item {} 
\sphinxAtStartPar
The low\sphinxhyphen{}pass component reduces noise and drift.

\item {} 
\sphinxAtStartPar
The high\sphinxhyphen{}pass component preserves responsiveness.

\end{itemize}


\section{Discrete implementation}
\label{\detokenize{appendix_p_complementary_filtering:discrete-implementation}}
\begin{sphinxVerbatim}[commandchars=\\\{\}]
\PYG{n}{filtered} \PYG{o}{=} \PYG{n}{alpha} \PYG{o}{*} \PYG{n}{new\PYGZus{}value} \PYG{o}{+} \PYG{p}{(}\PYG{l+m+mi}{1} \PYG{o}{\PYGZhy{}} \PYG{n}{alpha}\PYG{p}{)} \PYG{o}{*} \PYG{n}{previous\PYGZus{}filtered}
\end{sphinxVerbatim}


\section{Applied to derivative terms}
\label{\detokenize{appendix_p_complementary_filtering:applied-to-derivative-terms}}
\sphinxAtStartPar
Derivative signals are noisy. A complementary filter can smooth the derivative
term in a PD or PID controller, improving stability without a large delay.


\section{Applied to waypoint updates}
\label{\detokenize{appendix_p_complementary_filtering:applied-to-waypoint-updates}}
\sphinxAtStartPar
When target updates occur mid\sphinxhyphen{}motion, filtering avoids sudden jumps in desired
position and keeps the control output smooth.


\chapter{Appendix Q: CAN Bus}
\label{\detokenize{appendix_q_can_bus:appendix-q-can-bus}}\label{\detokenize{appendix_q_can_bus::doc}}

\section{Overview}
\label{\detokenize{appendix_q_can_bus:overview}}
\sphinxAtStartPar
Controller Area Network (CAN) is a robust bus protocol used in automotive and
robotic systems for reliable device communication.


\section{Frame structure}
\label{\detokenize{appendix_q_can_bus:frame-structure}}\begin{itemize}
\item {} 
\sphinxAtStartPar
Identifier: arbitration and message priority.

\item {} 
\sphinxAtStartPar
Data length: number of bytes (0 to 8 for classic CAN).

\item {} 
\sphinxAtStartPar
Data: payload bytes.

\item {} 
\sphinxAtStartPar
CRC: error detection.

\end{itemize}


\section{Arbitration}
\label{\detokenize{appendix_q_can_bus:arbitration}}
\sphinxAtStartPar
CAN uses non\sphinxhyphen{}destructive bitwise arbitration. The lowest identifier wins bus
access without corrupting messages.


\section{Bitrate and timing}
\label{\detokenize{appendix_q_can_bus:bitrate-and-timing}}
\sphinxAtStartPar
All nodes must share the same bitrate and timing parameters. Mismatched timing
causes errors and bus\sphinxhyphen{}off states.


\section{Error handling}
\label{\detokenize{appendix_q_can_bus:error-handling}}
\sphinxAtStartPar
CAN includes error counters and automatic retransmission. Persistent errors
can place a node into a bus\sphinxhyphen{}off state until recovery.


\section{Classic CAN vs CAN FD}
\label{\detokenize{appendix_q_can_bus:classic-can-vs-can-fd}}\begin{itemize}
\item {} 
\sphinxAtStartPar
Classic CAN: 8\sphinxhyphen{}byte payload limit.

\item {} 
\sphinxAtStartPar
CAN FD: larger payload and higher data phase bitrate.

\end{itemize}


\section{Applied to robotics}
\label{\detokenize{appendix_q_can_bus:applied-to-robotics}}
\sphinxAtStartPar
CAN is often used for motor controllers and sensor modules that require robust
communication in noisy electrical environments.


\chapter{Indices and tables}
\label{\detokenize{index:indices-and-tables}}\begin{itemize}
\item {} 
\sphinxAtStartPar
\DUrole{xref,std,std-ref}{genindex}

\item {} 
\sphinxAtStartPar
\DUrole{xref,std,std-ref}{modindex}

\item {} 
\sphinxAtStartPar
\DUrole{xref,std,std-ref}{search}

\end{itemize}



\renewcommand{\indexname}{Index}
\printindex
\end{document}